\documentclass[11pt]{article}
\usepackage[a4paper,margin=18mm]{geometry}
\usepackage{fontspec}
\setmainfont{Latin Modern Roman}
\usepackage{amsmath,amssymb,amsthm,amscd,graphicx,color}
\usepackage{xurl}
\usepackage[unicode,hidelinks]{hyperref}
\allowdisplaybreaks
\setlength\emergencystretch{3em}
\numberwithin{equation}{section}

\newtheorem{Theorem}{Theorem}[section]
\newtheorem{Lemma}[Theorem]{Lemma}
\newtheorem{Proposition}[Theorem]{Proposition}
 { \theoremstyle{definition}
\newtheorem{Definition}[Theorem]{Definition}
\newtheorem{Example}[Theorem]{Example}
\newtheorem{Remark}[Theorem]{Remark} }

\newcommand{\kl}{\langle}
\newcommand{\er}{\rangle}
\newcommand{\He}{\mathbf{D}}
\newcommand{\x}{\kl x \er}
\renewcommand{\r}{\kl r \er}
\newcommand{\ialpha}{i_{\alpha}}
\newcommand{\jbeta}{j_{\beta}}
\newcommand{\kgamma}{k_{\gamma}}
\newcommand{\jalpha}{j_{\alpha}}
\newcommand{\kbeta}{k_{\beta}}
\newcommand{\pgamma}{p_{\gamma}}
\newcommand{\malpha}{m_{\alpha}}
\newcommand{\mbeta}{m_{\beta}}
\newcommand{\mgamma}{m_{\gamma}}
\newcommand{\lbeta}{l_{\beta}}
\newcommand{\lalpha}{l_{\alpha}}
\newcommand{\palpha}{p_{\alpha}}
\newcommand{\kalpha}{k_{\alpha}}
\newcommand{\pbeta}{p_{\beta}}
\newcommand{\ngamma}{n_{\gamma}}
\newcommand{\igamma}{i_{\gamma}}
\newcommand{\jgamma}{j_{\gamma}}
\newcommand{\lgamma}{l_{\gamma}}
\newcommand{\nalpha}{n_{\alpha}}
\newcommand{\hgamma}{h_{\gamma}}
\newcommand{\halpha}{h_{\alpha}}
\newcommand{\idelta}{i_{\delta}}
\newcommand{\jdelta}{j_{\delta}}
\newcommand{\kepsilon}{k_{\epsilon}}
\newcommand{\rhoepsilon}{\rho_{\epsilon}}
\newcommand{\rhogamma}{\rho_{\gamma}}
\newcommand{\rhodelta}{\rho_{\delta}}
\newcommand{\ldelta}{l_{\delta}}
\newcommand{\pepsilon}{p_{\epsilon}}
\newcommand{\lepsilon}{l_{\epsilon}}
\newcommand{\ibeta}{i_{\beta}}
\newcommand{\iepsilon}{i_{\epsilon}}
\newcommand{\hepsilon}{h_{\epsilon}}
\newcommand{\hdelta}{h_{\delta}}
\newcommand{\pdelta}{p_{\delta}}



\makeatletter\newlabel{Perspectives}{{7}{}{Original source locator}{}{}}\makeatother
\makeatletter\newlabel{PainHelmSection}{{5}{}{Original source locator}{}{}}\makeatother
\makeatletter\newlabel{maintheorem2}{{1.5}{}{Original source locator}{}{}}\makeatother
\makeatletter\newlabel{maintheorem3}{{1.6}{}{Original source locator}{}{}}\makeatother
\makeatletter\newlabel{2}{{2}{}{Original source locator}{}{}}\makeatother
\begin{document}
\begin{center}\Large For Riemannian Painleve metrics with genuine variable blocks, the generalized Robertson conditions are equivalent to vanishing of every off-block-diagonal Ricci component\end{center}
\noindent\textbf{Stable ID:} 1703\par
\noindent\textbf{Authors:} Thierry Daude; Niky Kamran; Francois Nicoleau\par
\noindent\textbf{Original work:} Separability and Symmetry Operators for Painlev\&\#233; Metrics and their Conformal Deformations\par
\noindent\textbf{Version:} Official SIGMA 15 (2019) article 069, DOI 10.3842/SIGMA.2019.069; journal PDF and native TeX acquired 2026-09-30, exact bytes pinned by SHA256\par
\noindent\textbf{Selected task scope:} Exactly Theorem1.4(2) on a local Riemannian Painleve coordinate domain, restricted to 2<=r<n with at least one block dimension greater than one. S is an invertible generalized Staeckel matrix whose alpha-th row depends only on x\textasciicircum{}alpha; s\textasciicircum{}\{alpha1\} are its first-column cofactors and det(S)/s\textasciicircum{}\{alpha1\}>0. Each G\_alpha is a positive-definite block metric depending only on x\textasciicircum{}alpha, and g=sum\_alpha (det(S)/s\textasciicircum{}\{alpha1\})*G\_alpha. The generalized Robertson conditions (1.20), with the exact gamma\_\{i\_beta\} logarithmic determinant/cofactor expression, hold iff R\_\{j\_alpha,k\_beta\}=0 for every pair of distinct blocks and all their coordinate indices. No genericity assumption (3.5) is added, because the author supplies its removal in Section4. The product-separable PDE existence statement1.4(1) and subsequent commuting-operator theorem are not selected.\par
\noindent\textbf{Difficulty:} H2: The Staeckel one-variable identity does not control the additional curvature terms caused by blocks of dimension greater than one. The author calculates all cross-block Christoffel contractions and shows that their dependence on the arbitrary intrinsic block metrics cancels. Generalized Levi-Civita identities then remove the new remainder terms, giving the geometric Ricci characterization of grouped separation conditions.\par
\noindent\textbf{Editorial boundary note (not author text):} The full original geometric clause Theorem 1.4(2) is retained, rather than the neighboring separated-PDE existence clause. The selected task is a genuine grouped-variable Painleve metric. Complete author Staeckel/cofactor/Robertson notation, generalized Killing–Eisenhart and Levi-Civita identities, removal of the provisional distinct-eigenvalue hypothesis and the full Section 6.1 off-block Ricci calculation are included. Every summation block, intrinsic-metric cancellation and remaining Robertson mixed-log derivative is explicit author text. Product-separable PDE existence, symmetry operators and conformal R-separation are not counted; bounded author dummy-index slips are unchanged.\par
\noindent\textbf{Source:} \url{https://sigma-journal.com/2019/069/}\quad DOI: \url{https://doi.org/10.3842/SIGMA.2019.069}\par
\noindent\textbf{License and attribution:} Copyright retained by the named authors. Published by SIGMA under CC BY 4.0: \url{https://creativecommons.org/licenses/by/4.0/}. Source: \url{https://sigma-journal.com/about.html}. Adaptation consists of standalone excerpt formatting, cover and numbering restoration; original author language and mathematical text are retained.\par
\noindent\textbf{Editorial symbol notice (not author text):} This is an explicitly partial statement selection: the geometric equivalence1.4(2), not the neighbouring elliptic separated-solution existence assertion. The source has several bounded dummy-index slips in its long sums; the summation-block decomposition and cancelling remainder identities are given explicitly and retained unchanged. No standalone coefficient formula or edited curvature calculation is counted.\par
\noindent\textbf{External original locator Perspectives:} Original Section 7 discusses perspectives and open problems, mentioned only in commentary.\par
\noindent\textbf{External original locator PainHelmSection:} Original Section 5 proves the neighboring separated-PDE existence result, outside selected clause 1.4(2).\par
\noindent\textbf{External original locator maintheorem2:} Original Theorem 1.5 is the separate symmetry-operator assertion, mentioned only in the section overview.\par
\noindent\textbf{External original locator maintheorem3:} Original Theorem 1.6 is the separate conformal R-separation assertion, mentioned only in the section overview.\par
\noindent\textbf{External original locator 2:} Original Section 2 gives illustrative Painleve metrics, mentioned only for comparison.\par
\medskip\hrule\medskip\noindent\textbf{Author text begins below.}\par
\setcounter{section}{0}
\setcounter{subsection}{0}
\setcounter{subsubsection}{0}
\setcounter{Theorem}{0}
\setcounter{equation}{0}
% Original source lines 99--375
\section{Introduction and statement of results}
Painlev\'e metrics \cite{Pain1897, Per1990} are a class of Riemannian metrics that provide a broad generalization of the well-known separable metrics of St\"ackel \cite{Eis1934, Sta1897} to the case in which the Hamilton--Jacobi equation for the geodesic flow can be additively separated into partial differential equations depending on groups of independent variables rather ordinary differential equations resulting from a complete orthogonal separation. In particular, while St\"ackel metrics in dimension $n$ admit $n$ linearly independent Poisson-commuting quadratic first integrals of the geodesic flow, Painlev\'e metrics in dimension $n$ will admit only $r<n$ such integrals in general, where $r$ denotes the number of groups of variables.

Our goal in this paper is to carry out the extension to Painlev\'e metrics of the well-known results \cite{Ben2016, BCR1-2002,BCR2-2002, Eis1934, KaMi1980, KaMi1981, KaMi1984} which relate in the St\"ackel case the additive separation of variables for the Hamilton--Jacobi equation to the multiplicative separation of variables for the Helmholtz equation, and the existence of quadratic first integrals of the geodesic flow to that of symmetry operators for the Laplace--Beltrami operator. We shall thus obtain the generalization to Painlev\'e metrics of the Robertson separability conditions~\cite{Rob1928} for the Helmholtz equation for St\"ackel metrics, and a formulation of these generalized Robertson conditions in terms of the vanishing of the off-block diagonal components of the Ricci tensor, thereby extending the classical result proved by Eisenhart \cite{Eis1934} for St\"ackel metrics. We shall also show that when the generalized Robertson conditions are satisfied, there exist $r<n$ linearly independent second-order differential operators which commute between themselves and with the Laplace--Beltrami operator. These operators will be shown to admit the block-separable solutions of the Helmholtz equation as formal eigenfunctions, with the separation constants arising from the separation into groups variables as eigenvalues. Finally, we shall also discuss conformal deformations of Painlev\'e metrics satisfying a further generalization of the Robertson conditions, which is compatible with the separation into blocks of variables of the Helmholtz equation, leading to solutions which are $R$-separable in blocks.

Before describing our results in further detail, we should remark that independently of the interest of Painlev\'e metrics from the point of view of separation of variables, a key motivation for our study lies in the goal of constructing geometric models of manifolds with boundary endowed with Painlev\'e metrics, with the goal of investigating the anisotropic Calder\'on problem in this class of geometries. Recall that the anisotropic Calder\'on problem consists in recovering the metric of a Riemannian manifold with boundary from the knowledge of the Dirichlet-to-Neumann map defined by the Laplace--Beltrami operator. The anisotropic Calder\'on problem is at the center of a great amount of current research activity. We refer to \cite{DKN2016,DKN2019a, GT2013, KKL2001, KS2014, LU2001, LeU1989, Sa2013, Uhl2009, Uhl2014} and the references therein for surveys of recent results on this problem. We have recently investigated the anisotropic Calder\'on problem at fixed energy for geometric models consisting of classes of St\"ackel manifolds with boundary, where the separation of variables for the Helmholtz equation allows the decomposition of the Dirichlet-to-Neumann map onto a~basis of joint eigenfunctions of the symmetry operators resulting from the complete separation of variables, enabling us to obtain a series of uniqueness and non-uniqueness results for the Calder\'on problem, with no a-priori assumptions of analyticity, or on the existence of isometries \cite{DKN2015, DKN2016, DKN2017, DKN2018a, DKN2019b, Gob2016}. In the case of Painlev\'e metrics, the separation of the Helmholtz into groups of variables and the concomitant families of commuting symmetry operators admitted by these metrics will serve as an effective starting point for the investigation of the anisotropic Calder\'on problem in this more general setting.

In order to put the results of the present paper in context, we first briefly recall some well-known definitions and results pertaining to St\"ackel metrics and their separability properties. Throughout the paper, we shall assume for simplicity that the manifolds, metrics and maps being considered are smooth, although many of the results that we quote or obtain can be shown to hold under weaker differentiability properties. Recall \cite{BCR1-2002, Eis1934, KaMi1980, Sta1897} that a {\emph {St\"ackel metric}} on an $n$-dimensional manifold $M$ is a Riemannian metric $g$ for which there exist local coordinates $\big(x^{1},\dots,x^{n}\big)$ in which the metric has the expression
\begin{gather}\label{Stackelform}
{\rm d}s^{2}=g_{ij}{\rm d}x^{i}{\rm d}x^{j}=\frac{\det S}{s^{11}}\big({\rm d}x^{1}\big)^{2}+\cdots+\frac{\det S}{s^{n1}}\big({\rm d}x^{n}\big)^{2},
\end{gather}
where $S$ is a {\emph{St\"ackel matrix}}, that is a non-singular $n\times n$ matrix $S=(s_{ij})$ of the form
\begin{gather}\label{StackelMatrix}
S=\left(
\begin{matrix}
s_{11}\big(x^{1}\big)&\dots&s_{1n}\big(x^{1}\big) \\
\vdots& & \vdots \\
s_{n1}\big(x^{n}\big)&\dots&s_{nn}\big(x^{n}\big)
\end{matrix}
\right),
\end{gather}
and $s^{ij}$ denotes the cofactor of the component $s_{ij}$ of the matrix~$S$. St\"ackel matrices thus have the property that for each $1\leq i \leq n$, their $i$-th row depends only on the $i$-th local coordinate~$x^{i}$, and that the cofactor $s^{ij}$ is independent of the $i$-th local coordinate $x^{i}$. Furthermore, the diagonal components of the St\"ackel metric~(\ref{Stackelform}) are given by the inverses of the entries of the first row of the inverse St\"ackel matrix $A=S^{-1}$.

The importance of St\"ackel metrics stems from the fact that they constitute the most general class of Riemannian metrics admitting orthogonal coordinates for which the Hamilton--Jacobi equation
\begin{gather}\label{HJ}
g^{ij}\partial_{i}W\partial_{j}W=E,
\end{gather}
for the geodesic flow of $(M,g)$, where $E$ denotes a positive real constant, admits a complete integral obtained by additive separation of variables into ordinary differential equations. It is useful at this stage to recall that a complete integral of (\ref{HJ}) is defined as a parametrized family of solutions
\begin{gather}\label{complint}
W=W\big(x^{1},\dots,x^{n};a_{1},\dots,a_{n}\big),\qquad a_{1}:=E,
\end{gather}
defined over the domain $U\subset M$ of the local coordinates $\big(x^{1},\dots,x^{n}\big)$ and depending smoothly on $n$ parameters $(a_{1},\dots,a_{n})$ defined on an open subset $A\subset {\mathbb{R}}^{n}$, such that the rank condition
\begin{gather*}%\label{HJrank}
\det \left(\frac{\partial^{2}W}{\partial x^{i}\partial a_{j}}\right) \neq 0,
\end{gather*}
is satisfied throughout the open set $U\times A$.

It is easily verified that the Hamilton--Jacobi equation (\ref{HJ}) will admit an additively separable complete integral $W\big(x^{1},\dots, x^{n};a_{1},\dots a_{n}\big)$ of the form
\[
W=W_{1}\big(x^{1};a_{1},\dots, a_{n}\big)+\cdots+W_{n}\big(x^{n};a_{1},\dots, a_{n}\big),
\]
if and only if the summands $W_{\alpha}$ satisfy the set of separated ordinary differential equations given by
\[
\left(\frac{{\rm d}W_{i}}{{\rm d}x^{i}}\right)^{2}=\sum_{j=1}^{n}s_{ij}\big(x^{i}\big)a_{j}.
\]
The $n$ parameters $(a_{1},\dots , a_{n})$ appearing in the expression of the additively separable complete integral~(\ref{complint}) thus correspond to the separation constants arising from the complete separation of variables of the Hamilton--Jacobi equation into ordinary differential equations.
One of the key consequences of this complete separation of variables property is that the geodesic flow of an $n$-dimensional St\"ackel metric admits a linearly independent set of $n-1$ mutually Poisson-commuting quadratic first integrals, given by
\begin{gather}\label{FirstInt}
K_{(l)} = K_{(l)}^{ij} p_i p_j = \sum_{j=1}^{n}a_{lj} p_{j}^{2}, \qquad 2\leq l \leq n
\end{gather}
with
\[ \big\{K_{(l)},K_{(m)}\big\}=0,\qquad \text{for} \quad 1\leq l,m \leq n,
\]
where $A=(a_{ij})$ denotes as above the inverse of the St\"ackel matrix $S$ given by (\ref{StackelMatrix}). Note that with the notations of (\ref{FirstInt}), we have $K_{(1)}=H$, where
\[
H=g^{ij}p_{i}p_{j},
\]
denotes the Hamiltonian for the geodesic flow.

 A question closely related to the additive separation of the Hamilton--Jacobi equation is that of the complete multiplicative separation of the Helmholtz equation
 \begin{gather}\label{Helm}
-\Delta_{g} u =\lambda u,
\end{gather}
where
\[
\Delta_{g} = \frac{1}{\sqrt{|g|}}\partial_{i}\Big(\sqrt{|g|}g^{ij}\partial_{j}\Big),\qquad |g|=\det(g_{ij}),
\]
denotes the Laplace--Beltrami operator on~$(M,g)$ and $\lambda$ denotes a non vanishing real constant, into ordinary differential equations. By complete multiplicative separation, we mean, follo\-wing~\cite{BCR1-2002}, the existence of a parametrized family of solutions $u$ defined on a domain $U\subset M$ with local coordinates $\big(x^{1},\dots, x^{n}\big)$ of the form
\[
u\big(x^{1},\dots, x^{n};a_{1},\dots, a_{2n}\big)=\prod_{i=1}^{n}u_{i}\big(x^{1},\dots, x^{n};a_{1},\dots, a_{2n}\big),
\]
depending smoothly on $2n$ parameters $(a_{1},\dots, a_{2n})$ defined on open subset $A\subset {\mathbb R}^{2n}$, satisfying the rank condition
\begin{gather*}%\label{Helmrank}
\det \left(
\begin{matrix}
\dfrac{\partial v_{1}}{\partial a_{1}}&\dots&\dfrac{\partial v_{1}}{\partial a_{2n}} \\
\vdots& & \vdots \\
\dfrac{\partial v_{n}}{\partial a_{1}}&\dots&\dfrac {\partial v_{n}}{\partial a_{2n}}\vspace{1mm}\\
\dfrac{\partial w_{1}}{\partial a_{1}}&\dots&\dfrac{\partial w_{1}}{\partial a_{2n}} \\
\vdots& & \vdots \\
\dfrac{\partial w_{n}}{\partial a_{1}}&\dots&\dfrac {\partial w_{n}}{\partial a_{2n}}
\end{matrix}
\right) \neq 0,
\end{gather*}
at every point of $U\times A$, where
\[
v_{i}=\frac{u_{i}'}{u_{i}} \qquad w_{i}=\frac{u_{i}''}{u_{i}}.
\]
This separation requires that additional conditions, known as the {\emph{Robertson conditions}}, and given by
\[
\partial_{i}\Gamma_{j}=0,\qquad 1\leq i\neq j\leq n,
\]
where
\[
\Gamma_{i}=-\partial_{i}\left[\log \frac{(\det S)^{\frac{n}{2}-1}s^{i1}}{\big(s^{11}\cdots s^{n1}\big)^{\frac{1}{2}}}\right],
\]
be satisfied. We refer to \cite{BCR1-2002,BCR2-2002,Eis1934,KaMi1980,KaMi1981,Rob1928} for detailed proofs of the fact that the Robertson conditions are necessary and sufficient for complete multiplicative separation of the Helmholtz equation. The Robertson conditions were interpreted by Eisenhart \cite{Eis1934} in terms of the vanishing of the off-diagonal components of the Ricci tensor of the underlying St\"ackel metric, that is
\begin{gather}\label{classicalRobertson}
R_{ij}=0\qquad \text{for}\qquad 1\leq i\neq j\leq n.
\end{gather}
When the Robertson conditions are satisfied, the Poisson-commuting quadratic first integrals of the geodesic flow give rise to $n-1$ linearly independent second-order differential operators which commute with the Laplace--Beltrami operator and also commute pairwise. Rewriting the quadratic first integrals $K_{(l)}$ defined by (\ref{FirstInt}) in the form
\[
K_{(l)}=K_{(l)}^{ij}p_{i}p_{j},
\]
these commuting operators, denoted by $\Delta_{K_{(l)}}$, are of the form
\[
\Delta_{K_{(l)}}=\nabla_{i}\big(K_{(l)}^{ij}\nabla_{j}\big),\qquad 2\leq l \leq n,
\]
where $\nabla_i$ denotes the Levi-Civita connection on $(M,g)$. These operators, which are often referred to as {\emph{symmetry operators}}, admit the separable solutions of the Helmholtz equation as (formal) eigenfunctions. We will not give any further details on symmetry operators at this stage, nor shall we say anything about the proofs of the results we have just recalled since we shall shortly state and prove generalizations of these to the case of Painlev\'e metrics, which admit all St\"ackel metrics as a special case.

We conclude these preliminaries by remarking that the above setting may be expanded significantly by considering conformal deformations of St\"ackel metrics which are compatible with the complete separation of the Helmholtz equation into ordinary differential equations, thus giving rise to the more general notion of {\emph{$R$-separability}} for the Helmholtz equation. Again, we shall not give any additional details on these topics at this stage since conformal deformations and $R$-separability will be studied in the remainder of this paper in the more general setting of Painlev\'e metrics. We refer to \cite{BCR1-2002, BCR2-2002, BCR2005, CR2006,CR2007, KaMi1980, KaMi1981, KaMi1984} for lucid accounts of the key results on the separability and $R$-separability properties of St\"ackel metrics, their connection to Killing tensors, quadratic first integrals of the geodesic flow and symmetry operators for the Laplace--Beltrami operator. We also refer to \cite{Ben2016, Mi1988} for recent surveys on separability on Riemannian manifolds and to \cite{Car1977} for a penetrating analysis of the relations between quadratic first integrals of the geodesic flow, symmetry operators and conserved currents, in the general setting of Riemannian or pseudo-Riemannian manifolds.

With these preliminaries at hand, we are now ready to introduce the class of Painlev\'e met\-rics~\cite{Pain1897, Per1990}. As stated above, Painlev\'e metrics arise as a natural generalization of St\"ackel metrics to the case in which one no longer seeks complete separation of the Hamilton--Jacobi equation into ordinary differential equations, but rather separation into partial differential equations involving groups of variables. The separable coordinates admitted by Painlev\'e metrics are thus generally {\emph {not orthogonal}}, although they are orthogonal with respect to groups of variables. Let us recall that our goal in this paper is to carry out for Painlev\'e metrics the analogue of the separability and $R$-separability analyses of the Helmholtz equation which has been extensively worked out for St\"ackel metrics in \cite{BCR1-2002,BCR2-2002,CR2006,KaMi1980,KaMi1981, KaMi1984}, and to show that the separability into groups of variables gives rise to vector spaces of mutually commuting symmetry operators for the Laplace--Beltrami operator, the dimension of which is determined by the number of groups of variables. In particular, we will generalize to the case of Painlev\'e metrics the Robertson conditions and the characterization thereof in terms of the Ricci tensor. We now proceed to define the class of Painlev\'e metrics along lines similar to the ones used above for St\"ackel metrics.

Let $(M,g)$ be an $n$-dimensional Riemannian manifold and let ${\bf x}=\big(x^{1},\dots , x^{n}\big)$ denote a set of local coordinates on $M$. We shall consider partitions of ${\bf x}$ into $r\geq 2$ groups of local coordinates,
\[
{\bf x}=\big({{\bf x}^{1}},\dots,{{\bf x}^{r}}\big),
\]
where
\[
{\bf {x}}^{\alpha}=\big(x^{i_{\alpha}}\big),\qquad 1_{\alpha}\leq i_{\alpha} \leq l_{\alpha},\qquad \text{and}\qquad \sum_{\alpha=1}^{r}||l_{\alpha}||=n,
\]
where $||l_{\alpha}||$ denotes the number of local coordinates in the group of label $\alpha$. Latin indices $1\leq i,j,\ldots \leq n$ will be used to label the local coordinates on $M$, greek indices $\alpha,\beta,\dots$ to label the $r$ groups of local coordinates, and hybrid indices $i_{\alpha}$, $1_{\alpha} \leq i_{\alpha}\leq l_{\alpha}$ to denote the local coordinates within the group $\bf{x^{\alpha}}$. Unless there is an ambiguity in the notation being used, in which case we will write out the summation signs explicitly, we shall apply the summation convention with the above range of indices.
A {\em generalized St\"ackel matrix} is a non-singular $r\times r$ matrix-valued function $S$ on $M$ of the form
\begin{gather}\label{genStackelmat}
S=\left(
\begin{matrix}
s_{11}\big({{\bf x}^{1}}\big)&\dots&s_{1r}\big({{\bf x}^{1}}\big) \\
\vdots& & \vdots \\
s_{r1}\big({{\bf x}^{r}}\big)&\dots&s_{rr}\big({\bf x}^{r}\big)
\end{matrix}
\right),
\end{gather}
where for each $1\leq \alpha \leq r$,
\[
{\bf x}^{\alpha}=\big(x^{i_{\alpha}}\big),\qquad 1_{\alpha}\leq i_{\alpha} \leq l_{\alpha}.
\]
Let $s^{\alpha \beta}$ denote the cofactor of the component $s_{\alpha \beta}$ of $S$. We note that the cofactor $s^{\beta \gamma}$ will not depend on the group of variables ${\bf x}^{\beta}=\big(x^{i_{\beta}}\big)$, $1_{\beta}\leq i_{\beta} \leq l_{\beta}$. Moreover we shall assume that
\begin{gather} \label{StackelPositiveness}
\frac{\det S}{s^{\alpha 1}} > 0, \qquad \forall\, 1\leq \alpha \leq r ,
\end{gather}
in order to work with {\emph{Riemannian}} Painlev\'e metrics.

\begin{Definition} \label{Painleve}
Let $S$ be a generalized St\"ackel matrix satisfying (\ref{StackelPositiveness}). A Painlev\'e metric is a~Riemannian metric~$g$ for which there exist local coordinates ${\bf x}=\big({\bf x}^{1},\dots,{\bf x}^{r}\big)$ such that
\begin{gather}\label{Painleveform}
{\rm d}s^{2}=g_{ij}{\rm d}x^{i}{\rm d}x^{j}=\frac{\det S}{s^{11}} G_1 +\cdots+\frac{\det S}{s^{r1}} G_r,
\end{gather}
where each of the quadratic differential forms
\begin{gather}\label{blockmetric}
G_{\alpha}=G_{\alpha}({{\bf x}^{\alpha}})=\sum_{i_{\alpha}=1_{\alpha}}^{l_{\alpha}}\sum_{j_{\alpha}=1_{\alpha}}^{l_{\alpha}}G_{\alpha}\big({{\bf x}^{\alpha}}\big)_{i_{\alpha}j_{\alpha}}{\rm d}x^{i_{\alpha}}{\rm d}x^{j_{\alpha}},\qquad 1\leq \alpha \leq r,
\end{gather}
is positive-definite in its arguments and depends only on the group of variables ${\bf x}^{\alpha}$.
\end{Definition}
We may thus write the metric (\ref{Painleveform}) in block-diagonal form as
\[
{\rm d}s^{2}=\sum_{\alpha =1}^{r}\sum_{\ialpha = 1_{\alpha}}^{\lalpha}\sum_{\jalpha = 1_{\alpha}}^{\lalpha}(g_{\alpha})_{\ialpha \jalpha}{\rm d}x^{\ialpha}{\rm d}x^{\jalpha}=\sum_{\alpha =1}^{r} \frac{\det S}{s^{\alpha 1}} \sum_{\ialpha = 1_{\alpha}}^{\lalpha}\sum_{\jalpha = 1_{\alpha}}^{\lalpha}(G_{\alpha})_{i_{\alpha}j_{\alpha}}{\rm d}x^{i_{\alpha}}{\rm d}x^{j_{\alpha}}.
\]
It is important to note that even though Painlev\'e metrics (\ref{Painleveform}) are block-diagonal, and each quadratic differential form (\ref{blockmetric}) defines a Riemannian metric on the submanifolds defined by the level sets ${\bf x}^{\beta}={\bf{c}}^{\beta}$, $\beta\neq \alpha$, {\emph {a Painlev\'e metric is generally not a direct sum of Riemannian metrics, nor a warped product, except for special non-generic cases}}. We also note that Painlev\'e metrics of semi-Riemannian (and in particular Lorentzian) signature can readily be defined by modifying the requirement that each of the quadratic differential forms $G_{\alpha}$ given by (\ref{blockmetric}) be positive-definite to one in which $G_{{\alpha}}$ is assumed to be of signature $(p_{\alpha},q_{\alpha})$ with $p_{\alpha}+q_{\alpha}=l_{\alpha}$. Finally we also remark that the Painlev\'e form (\ref{Painleveform}) is obviously invariant under smooth and invertible changes of coordinates of the form ${\tilde{\bf x}}^{\alpha}={\bf{f}}^{\alpha}({\bf x}^{\alpha})$, where $1\leq \alpha \leq r$.

Let us call \emph{block orthogonal coordinates} a system of coordinates $({\bf x}^{\alpha})$ such that the metric $g$ has the form
\begin{gather*} %\label{BlockOrthoCoord}
 g=\sum_{\alpha =1}^{r} c_\alpha G_\alpha = \sum_{\alpha =1}^{r} c_\alpha \sum_{\ialpha = 1_{\alpha}}^{\lalpha}\sum_{\jalpha = 1_{\alpha}}^{\lalpha}(G_{\alpha})_{i_{\alpha}j_{\alpha}}{\rm d}x^{i_{\alpha}}{\rm d}x^{j_{\alpha}},
\end{gather*}
where $c_\alpha$ are non-vanishing scalar functions on $M$ and the metrics $G_\alpha$ are given by (\ref{blockmetric}).
In analogy with the St\"ackel case, we have

\begin{Proposition} \label{HJSeparability-Painleve}A metric $g$ is of the Painlev\'e form \eqref{Painleveform} if and only if there exist block orthogonal coordinates such that the Hamilton--Jacobi equation
\[
g^{ij}\partial_{i}W\partial_{j}W=E,
\]
admits a parametrized family of solutions which is sum-separable into groups of variables, of the form
\begin{gather}\label{Painsepform}
W=W_{1}\big({\bf {x}}^{1};{{a}}_{1},\dots ,{{a}}_{r}\big)+\cdots+W_{r}\big({\bf {x}}^{r};{{{a}}_{1},\dots, {a}}_{r}\big),
\end{gather}
depending smoothly on $r$ arbitrary real constants $({{a}}_{1}:=E,a_{2},\dots ,{{a}}_{r})$ defined on an open subset $A\subset \mathbb{R}^{r}$, and satisfying the rank condition
\begin{gather}\label{rankPainHJ}
\det \left(\sum_{\ialpha =1_{\alpha}}^{\lalpha}\sum_{\jalpha =1_{\alpha}}^{\lalpha}\big(G^{\alpha}\big)^{\ialpha \jalpha}\left(\frac{\partial W_{\alpha}}{\partial x^{\ialpha}}\right)\left( \frac{\partial^{2}W_{\alpha}}{\partial a_{\gamma}\partial x^{\jalpha}}\right)\right)\neq 0,
\end{gather}
where
\[
G^{\alpha} = (G_{\alpha})^{-1 }.
\]
\end{Proposition}
This Proposition will be proved in Section \ref{PainHJSection} as well as other (intrinsic) characterizations of Painlev\'e metrics.

In further analogy with the St\"ackel case, we now recall that Painlev\'e metrics admit $r$ linearly independent quadratic first integrals of the geodesic flow which are Poisson commuting. Indeed, the summands $W_{\alpha}$ appearing in (\ref{Painsepform}) satisfy the following set of first-order PDEs \cite{Per1990}
\begin{gather}
{\mathcal{F}}_{1}\left({\bf {x}}^{1},\frac{\partial W_{1}}{\partial {\bf x}^{1}}\right) =a_{1}s_{11}\big({\bf {x}}^{1}\big)+\cdots +a_{r}s_{r1}\big({\bf {x}}^{1}\big),\nonumber\\ \quad \vdots \nonumber\\ {\mathcal{F}}_{r}\left({\bf {x}}^{r},\frac{\partial W_{r}}{\partial {\bf x}^{r}}\right) =a_{1}s_{r1}\big({\bf {x}}^{r}\big)+\cdots +a_{r}s_{rr}\big({\bf {x}}^{r}\big),\label{HJsep}
\end{gather}
where
\begin{gather}\label{F}
{\mathcal{F}}_{\alpha}=\big(G^{\alpha}\big)^{i_{\alpha}j_{\alpha}}\big({\bf {x}}^{\alpha}\big)p_{i_{\alpha}}p_{j_{\alpha}},\qquad 1\leq \alpha \leq r,
\end{gather}
and where the $(a_{\alpha})$ are arbitrary real separation constants.
It follows now directly from the separated equations (\ref{HJsep}) and from the fact that the generalized St\"ackel matrix $S$ is non-singular that
one obtains $r$ linearly independent Poisson-commuting quadratic first integrals~$K_{(\alpha)}$ of the geodesic flow by solving for the $r$ separation constants $(a_{\alpha})$ from the separated equations (\ref{HJsep}). These are explicitly given by (see~\cite{Per1990})
\[
K_{(\alpha)}=\sum_{\beta =1}^{r}\big(S^{-1}\big)_{\alpha \beta}\mathcal{F}_{\beta},\qquad 1\leq \beta \leq r,
\]
or equivalently
\[
K_{(\alpha)}=K_{(\alpha)}^{ij}p_{i}p_{j},
\]
where for each $1\leq \alpha \leq r$, $\big(K_{(\alpha)}^{ij}\big)$ is a symmetric block-diagonal tensor defined by
\begin{gather}\label{Killingtensor}
K_{(\alpha)}^{i_{\beta}j_{\beta}}=\frac{s^{\beta \alpha}}{\det S}\big(G^{\beta}\big)^{i_{\beta}j_{\beta}},\qquad K_{(\alpha)}^{i_{\beta}j_{\gamma}}=0\qquad \text{for all}\qquad 1 \leq \beta \neq \gamma \leq r.
\end{gather}
These quadratic first integrals satisfy
\begin{gather}\label{PoissonCommute}
\big\{K_{(\alpha)}, K_{(\beta)}\big\}=0, \qquad 1\leq \alpha,\beta \leq r, \qquad \text{where}\qquad K_{(1)}=H.
\end{gather}
The condition $\{K_{(\alpha)},H\}=0$ is equivalent to $\big(K_{(\alpha)}^{ij}\big)$ being a symmetric {\emph{Killing tensor}}, that is
\[
\nabla_{i}K_{(\alpha) jk}+\nabla_{j}K_{(\alpha) ki}+\nabla_{k}K_{(\alpha) ij}=0.
\]
The commutation relations (\ref{PoissonCommute}) are thus equivalent to the vanishing of the Schouten brackets of the pairs of Killing tensors $(K_{(\alpha) ij}),(K_{(\beta) ij})$.

There exist a few classical examples of Painlev\'e metrics in the litterature. They include for instance the {\emph{di Pirro metrics}} \cite{diPirro1896, Per1990}, for which the Hamiltonian of the geodesic flow is of the form
\begin{gather}\label{diPirro}
H=g^{ij}p_{i}p_{j}=\big(c_{12}\big(x^{1},x^{2}\big)+c_{3}\big(x^{3}\big)\big)^{-1}\big[a_{1}\big(x^{1},x^{2}\big)p_{1}^{2}+a_{2}\big(x^{1},x^{2}\big)p_{2}^{2} +a_{3}\big(x^{3}\big)p_{3}^{2}\big].
\end{gather}
It may indeed be verified directly that the function
\[
K=\big(c_{12}\big(x^{1},x^{2}\big)+c_{3}\big(x^{3}\big)\big)^{-1}\big[c_{3}\big(x^{3}\big)(a_{1}\big(x^{1},x^{2}\big)p_{1}^{2} +a_{2}\big(x^{1},x^{2}\big)p_{2}^{2})-c_{12}\big(x^{1},x^{2}\big)a_{3}\big(x^{3}\big)p_{3}^{2}\big].
\]
Poisson-commutes with $H$, and thus defines a Killing tensor, which together with the metric tensor generates a maximal linearly independent set of Killing tensors for generic choices of the metric functions $c_{12}$, $a_{1}$, $a_{2}$, $a_{3}$, $c_{3}$ in~(\ref{diPirro}). Painlev\'e metrics also appear in the context of geodesically equivalent metrics as metrics admitting projective symmetries, see \cite{Top1999, Top2002}, and also as instances of 4-dimensional Lorentzian metrics admitting a Killing tensor \cite{HaMa1978} (see Section~\ref{Perspectives} for further remarks on the latter point). Finally, we mention the recent paper by Chanu and Rastelli~\cite{CR2019} that provides a classification of Painlev\'e metrics with vanishing Riemann tensor in dimension~$3$,  i.e., in~$\mathbb{E}_3$. We will give some examples of Painlev\'e metrics in all dimensions satisfying the generalized Robertson conditions (see below) as well as a catalogue of such metrics in dimensions $2$, $3$, $4$ in Section \ref{2}.

As we stated above, our main goal in this paper is to investigate for the class of Painlev\'e metrics the closely related question of product separability for the Helmholtz equation (\ref{Helm}), and the relationship between quadratic first integrals of the geodesic flow and symmetry operators for the Laplace--Beltrami equation. The Laplace--Beltrami operator for a Painlev\'e metric $g$ given by~(\ref{Painleveform}) can be expressed in terms of the generalized St\"ackel matrix $S$ and the Laplace--Beltrami operators for the $r$ Riemannian metrics $G_{\beta}$, $1\leq \beta \leq r$, defined by~(\ref{blockmetric}), corresponding to the blocks of variables ${\bf x}^{\beta}$, $1\leq \beta \leq r$. We have
\begin{gather}
\Delta_{g} u =\sum_{\beta=1}^{r}\Bigg[\left(\frac{s^{\beta 1}}{\det S}\right)\Bigg\{\Delta_{G_{\beta}}u\nonumber\\
\hphantom{\Delta_{g} u =}{}+\sum_{i_{\beta}=1_{\beta}}^{l_{\beta}}\sum_{{j}_{\beta}=1_{\beta}}^{l_{\beta}}\big(G^{\beta}\big)^{i_{\beta} j_{\beta}}\left[\partial_{i_{\beta}}\left(\log \frac{(\det S)^{\frac{n}{2}-1}s^{\beta 1}}{\big(s^{11}\big)^{\frac{l_1}{2}} \cdots \big(s^{r1}\big)^{\frac{l_r}{2}}}\right)\right]\partial_{j_{\beta}}u\Bigg\}\Bigg],\label{Laplacian}
\end{gather}
where $\Delta_{G_{\beta}}$ denotes the Laplace--Beltrami operator for the Riemannian metric $G_{\beta}$, that is
\begin{gather*}%\label{blockLaplacian}
\Delta_{G_{\beta}} = \sum_{i_{\beta}=1}^{\lbeta}\sum_{j_{\beta}=1}^{\lbeta}\frac{1}{\sqrt{|G_{\beta}|}}\partial_{i_{\beta}}\Big(\sqrt{|G_{\beta}|}
\big(G^{\beta}\big)^{i_{\beta}j_{\beta}}\partial_{j_{\beta}}\Big),\qquad |G_{\beta}|=\det((G_{\beta})_{i_{\beta}j_{\beta}}).
\end{gather*}

We now state our main results, the proofs of which will be given in Section \ref{ProofsofMainTheorems}. We first define the \emph{generalized Robertson conditions}, in analogy with the classical Robertson conditions (\ref{classicalRobertson}) for St\"ackel metrics.

\begin{Definition} \label{defiRC}
A Painlev\'e metric $g$ is said to satisfy the generalized Robertson conditions if and only if the differential conditions
	\begin{gather}\label{GenRobertson}
\partial_{j_{\alpha}}\gamma_{i_{\beta}}=0,\quad \textrm{for all} \quad 1\leq \alpha \neq \beta \leq r, \quad 1_{\alpha}\leq i_{\alpha} \leq l_{\alpha},\quad 1_{\beta}\leq i_{\beta} \leq l_{\beta},
\end{gather}
where
\[
\gamma_{i_{\beta}}=-\partial_{i_{\beta}}\left[\log \frac{(\det S)^{\frac{n}{2}-1}s^{\beta 1}}{\big(s^{11}\big)^{\frac{l_1}{2}} \cdots \big(s^{r1}\big)^{\frac{l_r}{2}}}\right],
\]
are satisfied.
\end{Definition}
The generalized Robertson conditions (\ref{GenRobertson}) imply that
\begin{gather}\label{GenRobertsonsimpl}
\partial_{\ialpha}\gamma^{\jbeta}=0, \quad \textrm{for all} \quad 1\leq \alpha \neq \beta \leq r,\\
\label{gammaup}
\gamma^{j_{\beta}}:= \sum_{i_{\beta}=1_{\beta}}^{l_{\beta}}\big(G^{\beta}\big)^{i_{\beta}j_{\beta}} \gamma_{i_{\beta}}.
\end{gather}
We shall be working with both the forms (\ref{GenRobertson}) and (\ref{GenRobertsonsimpl}) of these conditions. Note that if the Robertson conditions hold, then the positive Laplace--Beltrami operator can be written in a synthetic form as
\begin{gather*}%\label{LaplacianSimp}
 -\Delta_{g} u  =   \sum_{\beta=1}^{r} \left(\frac{s^{\beta 1}}{\det S}\right)\left\{-\Delta_{G_{\beta}}u + \sum_{{j}_{\beta}=1_{\beta}}^{l_{\beta}} \gamma^{j_{\beta}} \partial_{j_{\beta}}u\right\}
   =   \sum_{\beta=1}^{r} \left(\frac{s^{\beta 1}}{\det S}\right) B_\beta u , \nonumber
\end{gather*}
where the operators $B_\beta = -\Delta_{G_{\beta}} + \sum\limits_{{j}_{\beta}=1_{\beta}}^{l_{\beta}} \gamma^{j_{\beta}} \partial_{j_{\beta}}$ only depend on the group of variables $\textbf{x}^\beta$.

As will be seen in Section \ref{PainHelmSection}, the generalization of the notion of complete multiplicative separation for the Helmholtz equation to the case of separation in terms of groups of variables is given by considering a parametrized family of product-separable solutions of the form
\begin{gather} \label{SeparabilitySchro}
u=\prod_{\beta =1}^{r}u_{\beta}\big({\bf x}^{\beta};a_{1},\dots,a_{r}\big),
\end{gather}
satisfying the rank condition
\begin{gather} \label{RankPainSchro}
 \det  \left( \partial_{a_\alpha} \left( \frac{B_\beta u_\beta}{u_\beta} \right) \right) \ne 0,
\end{gather}
where we assume that $u_{\beta}\neq 0$.

Our first result states the separability conditions for the Helmholtz equation, and gives their interpretation in terms of the vanishing of the off-block diagonal components of the Ricci tensor:


\setcounter{section}{1}
\setcounter{subsection}{0}
\setcounter{subsubsection}{0}
\setcounter{Theorem}{3}
\setcounter{equation}{24}
% Original source lines 376--377
\begin{Theorem}\label{maintheorem1}\quad
\begin{enumerate}\itemsep=0pt

\setcounter{section}{1}
\setcounter{subsection}{0}
\setcounter{subsubsection}{0}
\setcounter{Theorem}{4}
\setcounter{equation}{26}
% Original source lines 388--393
\item[$2)$] The conditions \eqref{GenRobertson} may be written equivalently in terms of the Ricci tensor of the Painlev\'e metric \eqref{Painleveform} as
\begin{gather}\label{Ricciform}
R_{j_{\alpha}k_{\beta}}=0, \quad {\mbox{for all}} \quad 1\leq \alpha \neq \beta \leq r, \quad {\mbox{and}} \quad 1_{\alpha}\leq j_{\alpha} \leq l_{\alpha}, \quad 1_{\beta}\leq k_{\beta}\leq l_{\beta}.
\end{gather}
\end{enumerate}
\end{Theorem}

\setcounter{section}{2}
\setcounter{subsection}{0}
\setcounter{subsubsection}{0}
\setcounter{Theorem}{5}
\setcounter{equation}{5}
% Original source lines 744--1027
\section{Generalized Killing--Eisenhart and Levi-Civita conditions}

The proofs of the main results of our paper, that is Theorems \ref{maintheorem1}, \ref{maintheorem2} and \ref{maintheorem3}, make use of generalizations to Painlev\'e metrics of the classical Killing--Eisenhart equations and Levi-Civita separability conditions which hold for St\"ackel metrics (see for example \cite{BCR1-2002, BCR2-2002, Eis1934, KaMi1980, KaMi1981, Per1990}). We present these generalizations in the form of the following two lemmas, beginning with the Killing--Eisenhart equations. Thus in analogy with the St\"ackel case, we introduce the quantities
\begin{gather}\label{defrho}
\rho_{\beta \gamma}:=\frac{s^{\gamma \beta}}{s^{\gamma 1}}.
\end{gather}
Note that by the assumption (\ref{StackelPositiveness}), we have $s^{\gamma 1}\neq 0$.
The following lemma gives the generalization to the case of Painlev\'e metrics of the Killing--Eisenhart equations given in \cite{BCR1-2002,BCR2-2002,Eis1934,KaMi1980,KaMi1981} for St\"ackel metrics:
\begin{Lemma} \label{KE}
We have, for all $1\leq \beta, \delta, \gamma \leq r$, the identities
\begin{gather}\label{genKE}
\partial_{j_{\gamma}}(\rho_{ \beta \delta})= (\rho_{\beta \gamma}-\rho_{\beta \delta })\left(\partial_{j_{\gamma}}{\log \frac{s^{\delta 1}}{\det S}}\right).
\end{gather}
\end{Lemma}
We will refer to \eqref{genKE} as the {\emph{generalized Killing--Eisenhart equations.}}
\begin{proof}The Poisson bracket relations (\ref{PoissonCommute}) imply
\begin{gather}\label{Painlevecomm}
\sum_{\pgamma=1}^{\lgamma}\big(\big(\partial_{\pgamma} K_{(1)}^{\idelta \jdelta}\big) K_{(\beta)}^{\pgamma \kgamma}-\big(\partial_{\pgamma} K_{(\beta)}^{\idelta \jdelta}\big) K_{(1)}^{\pgamma \kgamma}\big)=0,
\end{gather}
where $1\leq \idelta, \jdelta \leq \ldelta$, $1\leq \kgamma \leq \lgamma$ and $1\leq \delta, \gamma \leq r$. Using the expressions~(\ref{Killingtensor}) of the Killing tensors $\big(K_{(\beta)}^{\idelta \jdelta}\big)$, the fact that $\partial_{\igamma}\big(G^{\beta}\big)^{\ibeta \jbeta}=0$ for $\beta \neq \gamma$, and the fact that each of the $\lgamma \times \lgamma$ matrices $\big(\big(G^{\gamma}\big)^{\igamma \jgamma}\big)$, is invertible, we obtain that the relations~(\ref{Painlevecomm}) are equivalent to
\begin{gather}\label{Stackelcomm}
\left[\partial_{\pgamma}\left(\frac{s^{\delta 1}}{\det S}\right)\right] \frac{s^{\gamma \beta}}{\det S}-\left[\partial_{\pgamma}\left(\frac{s^{\delta \beta}}{\det S}\right)\right] \frac{s^{\gamma 1}}{\det S}=0,
\end{gather}
where $1\leq \delta, \gamma \leq r$, $1_{\gamma}\leq \pgamma \leq \lgamma $. In particular, the relations (\ref{Painlevecomm}) are independent of the block metrics (\ref{blockmetric}). Setting $\delta = \gamma$ in~(\ref{Stackelcomm}), we obtain
\[
\partial_{\pgamma}\left(\frac{s^{\gamma \beta}}{s^{\gamma 1}}\right)=0,
\]
for $1_{\gamma}\leq \pgamma \leq \lgamma$, so that using the definition (\ref{defrho}) of the quantities $\rho_{\beta \gamma}$, the relations (\ref{Stackelcomm}) take the form
\[
\partial_{\pgamma}\left(\rho_{ \beta \delta}\frac{s^{\delta 1}}{\det S}\right)\frac{s^{\gamma 1}}{\det S}=\left(\partial_{\pgamma}\left(\frac{s^{\delta 1}}{\det S}\right)\right)\rho_{\beta \gamma} \frac{s^{\gamma 1}}{\det S},
\]
which in turn reduces to (\ref{genKE}), thus proving our claim.
\end{proof}

In order to state the generalization to Painlev\'e metrics of the classical Levi-Civita conditions which hold for St\"ackel metrics, we make the hypothesis
\begin{gather}\label{H}
\rho_{\beta \delta}\neq \rho_{\beta \epsilon},\qquad \forall\, 1\leq \beta \leq r,\quad \forall\, 1\leq \delta\neq \epsilon \leq r.
\end{gather}

The generalization of the Levi-Civita conditions to the Painlev\'e case is now given by the following:
\begin{Lemma} \label{LeviCivita}
A generalized St\"ackel matrix \eqref{genStackelmat} corresponding to a Painlev\'e metric~\eqref{Painleveform} for which the genericity hypothesis~\eqref{H} holds true satisfies, for all $1\leq \beta, \gamma \leq r$,
the identities
\begin{gather}
\left(\partial_{\jalpha}\log\left(\frac{s^{\gamma 1}}{\det S}\right)\right) \left(\partial_{\kbeta}\log\left(\frac{s^{\alpha 1}}{\det S}\right)\right)+\left(\partial_{\jalpha}\log\left(\frac{s^{\beta 1}}{\det S}\right)\right)\left(\partial_{\kbeta}\log\left(\frac{s^{\gamma 1}}{\det S}\right)\right)\nonumber\\
\qquad{} -\frac{\det S}{s^{\gamma 1}}\partial_{\jalpha}\partial_{\kbeta}\left(\frac{s^{\gamma 1}}{\det S}\right)=0.\label{genLeviCivita}
\end{gather}
In particular, we have the identities
\begin{gather}\label{secondorder}
\frac{\partial}{\partial x^{\jalpha}}\frac{\partial }{\partial x^{\kbeta}}\left(\frac{\det S}{s^{\alpha1}s^{\beta1}}\right)=0,
\end{gather}
for all
$1\leq \alpha, \beta \leq r$.
\end{Lemma}
We will likewise refer to the identities (\ref{genLeviCivita}) as the {\emph{generalized Levi-Civita conditions}}.
\begin{Remark}\quad
\begin{enumerate}\itemsep=0pt
\item[1)] In the case of St\"ackel metrics, that is when $r=n$, the conditions (\ref{genLeviCivita}) reduce to the classical Levi-Civita conditions, given by
\begin{gather*}
\left(\partial_{j}\log\left(\frac{s^{l1}}{\det S}\right)\right)\left(\partial_{k}\log\left(\frac{s^{j1}}{\det S}\right)\right)+\left(\partial_{j}\log\left(\frac{s^{k1}}{\det S}\right)\right)\left(\partial_{k}\log \left(\frac{s^{l1}}{\det S}\right)\right)\\
\qquad{} -\frac{\det S}{s^{l1}}\partial_{j}\partial_{k}\left(\frac{s^{l1}}{\det S}\right)=0.
\end{gather*}
\item[2)] We shall show at the end of the next Section~\ref{PainHJSection} that the generalized Levi-Civita condi\-tions~(\ref{genLeviCivita}) hold in fact for all Painlev\'e metrics~(\ref{Painleveform}) without assuming our genericity hypothesis~(\ref{H}). Nevertheless, it is easier to obtain them from the Killing--Eisenhart equations under the assumption~(\ref{H}) as we do below.
\end{enumerate}
\end{Remark}
\begin{proof}
The general idea behind the proof is similar to the one that is used in the classical St\"ackel case, and is based on expressing the integrability conditions for the generalized Killing--Eisenhart (\ref{genKE}), with additional twist resulting from the fact that the separation is in groups of variables only. We let $1\leq \alpha \neq \beta \leq r$ denote fixed indices and introduce the simplified notation
\[
\rho_{\delta}:=\rho_{ \beta \delta}=\frac{s^{\delta \beta}}{s^{\delta 1}},
\]
so as to make the expressions a bit more compact. The generalized Killing--Eisenhart equa\-tions~(\ref{genKE}) thus take the form
\begin{gather*}%\label{simpgenKE}
\partial_{\pgamma} \rho_{\delta}= (\rho_{\gamma}-\rho_{\delta })\left(\partial_{\pgamma}\left(\log \frac{s^{\delta 1}}{\det S}\right)\right),
\end{gather*}
where $1\leq \delta,\gamma \leq r$. The integrability conditions
\begin{gather}\label{intKE}
\partial_{\kepsilon}\big(\partial_{\pgamma} \rhodelta \big)=\partial_{\pgamma}\big(\partial_{\kepsilon} \rhodelta \big)
\end{gather}
are now easily obtained. We have
\begin{gather*}
\partial_{\kepsilon} (\partial_{\pgamma} \rhodelta )= \left[(\rhoepsilon-\rhogamma)\partial_{\kepsilon}\log\left(\frac{s^{\gamma 1}}{\det S}\right)-(\rhoepsilon-\rhodelta)\partial_{\kepsilon}\log\left(\frac{s^{\delta 1}}{\det S}\right)\right] \partial_{\pgamma} \log \frac{s^{\delta 1}}{\det S} \\
\hphantom{\partial_{\kepsilon} (\partial_{\pgamma} \rhodelta )=}{} + (\rhogamma-\rhodelta)\partial_{\kepsilon}\partial_{\pgamma}\log\left(\frac{s^{\delta 1}}{\det S}\right),
\end{gather*}
so that (\ref{intKE}) becomes
\begin{gather}
(\rho_{\gamma}-\rho_{\epsilon})\left[\left(\partial_{\kepsilon}\log\left(\frac{s^{\gamma 1}}{\det S}\right)\right)\left(\partial_{\pgamma}\log\left(\frac{s^{\delta 1}}{\det S}\right)\right)\right.\nonumber\\
\qquad{} +\left(\partial_{\kepsilon}\log\left(\frac{s^{\delta 1}}{\det S}\right)\right)\left(\partial_{\pgamma}\log\left(\frac{s^{\epsilon 1}}{\det S}\right)\right) \nonumber\\
\left.\qquad{} -\left(\partial_{\kepsilon}\log\left(\frac{s^{\delta 1}}{\det S}\right)\right)\left(\partial_{\pgamma}\log\left(\frac{s^{\delta 1}}{\det S}\right)\right)-\partial_{\kepsilon}\partial_{\pgamma}\log\left(\frac{s^{\delta 1}}{\det S}\right)\right]=0,\label{compat}
\end{gather}
where $1\leq \epsilon \neq \gamma \leq r$ and $1\leq \delta \leq r$. Using the rank hypothesis (\ref{H}), expanding the logarithmic second derivative in (\ref{compat}) using the identity
\begin{gather}\label{identity}
\partial_{xy}\log f = \frac{1}{f}\partial_{x}\partial_{y}f-(\partial_{x}\log f)(\partial_{y}\log f)
\end{gather}
and
 relabeling the indices, we obtain
\begin{gather*}
\left(\partial_{\jalpha}\log\left(\frac{s^{\gamma 1}}{\det S}\right)\right) \left(\partial_{\kbeta}\log\left(\frac{s^{\alpha 1}}{\det S}\right)\right)+\left(\partial_{\jalpha}\log\left(\frac{s^{\beta 1}}{\det S}\right)\right)\left(\partial_{\kbeta}\log\left(\frac{s^{\gamma 1}}{\det S}\right)\right)\\
\qquad{}-\frac{\det S}{s^{\gamma 1}}\partial_{\jalpha}\partial_{\kbeta}\left(\frac{s^{\gamma 1}}{\det S}\right)=0,
\end{gather*}
which is precisely (\ref{genLeviCivita}). Finally, the relations~(\ref{secondorder}) are obtained by setting $\delta = \epsilon$ in the integrability condition (\ref{compat}), using the identity (\ref{identity}), and the fact that the cofactors $s^{\gamma \alpha}$ and $s^{\epsilon \alpha}$ are independent of the groups of variables ${\bf x}^{\gamma}$ and ${\bf x}^{\epsilon}$ respectively.
\end{proof}

\section{Different characterizations of Painlev\'e metrics}\label{PainHJSection}

Let us start with the characterization of Painlev\'e metrics in terms of complete additive separability of the Hamilton--Jacobi equations stated in Proposition \ref{HJSeparability-Painleve}.

\begin{proof}[Proof of Proposition \ref{HJSeparability-Painleve}]

In block orthogonal coordinates $(\textbf{x}^\alpha)$ associated to the metric
\[
 g = \sum_{\alpha = 1}^r c_\alpha G_\alpha,
\]
the Hamilton--Jacobi equation (\ref{HJ}) reads
\begin{gather} \label{HJ1}
 \sum_{\alpha =1}^{r} c^\alpha \sum_{\ialpha =1_{\alpha}}^{\lalpha}\sum_{\jalpha =1_{\alpha}}^{\lalpha} \big(G^{\alpha}\big)^{\ialpha \jalpha}\frac{\partial W}{\partial x^{\ialpha} }\frac{\partial W}{\partial x^{\jalpha}}=E=:a_{1},
\end{gather}
where $c^\alpha = (c_\alpha)^{-1}$. Assume that there exists a solution $W$ in the block-separable form~(\ref{Painsepform}) satisfying the completeness condition~(\ref{rankPainHJ}). Then (\ref{HJ1}) can be written as
\begin{gather} \label{HJ2}
 \sum_{\alpha =1}^{r} c^\alpha \sum_{\ialpha =1_{\alpha}}^{\lalpha}\sum_{\jalpha =1_{\alpha}}^{\lalpha} \big(G^{\alpha}\big)^{\ialpha \jalpha}\frac{\partial W_\alpha}{\partial x^{\ialpha} }\frac{\partial W_\alpha}{\partial x^{\jalpha}}=E=:a_{1}.
\end{gather}
Differentiating (\ref{HJ2}) with respect to $a_\beta$, we get
\begin{gather} \label{HJ3}
 \sum_{\alpha =1}^{r} c^\alpha \sum_{\ialpha =1_{\alpha}}^{\lalpha}\sum_{\jalpha =1_{\alpha}}^{\lalpha} 2 \big(G^{\alpha}\big)^{\ialpha \jalpha}\frac{\partial W_\alpha}{\partial x^{\ialpha} } \frac{\partial^2 W_\alpha}{\partial a_\beta \partial x^{\jalpha}} = \delta_{1\beta} .
\end{gather}
From (\ref{Painsepform}) and (\ref{rankPainHJ}), it follows that the family of matrices
\begin{gather*} %\label{SM}
 S(a_1,\dots,a_r) = (S_{\alpha \beta})(a_1,\dots,a_r) := 2 \big(G^{\alpha}\big)^{\ialpha \jalpha}\frac{\partial W_\alpha}{\partial x^{\ialpha} } \frac{\partial^2 W_\alpha}{\partial a_\beta \partial x^{\jalpha}}
\end{gather*}
are non-singular st\"ackel matrices of rank $r$ for all $(a_1,\dots,a_r) \in U$. Using the invertibility of $S(a_1,\dots,a_r)$ which is equivalent to the completeness condition (\ref{rankPainHJ}), we get immediately from~(\ref{HJ3}) that
\[
 c^\alpha = \frac{s^{\alpha 1}}{\det S},
\]
which proves that $g$ is a Painlev\'e metric.

Conversely, assume that $g$ is a Painlev\'e metric of the form (\ref{Painleveform}). Then the Hamilton--Jacobi equation (\ref{HJ}) takes the form
\begin{gather}\label{HJPainform}
\sum_{\alpha =1}^{r} \left(\frac{s^{\alpha 1}}{\det S}\right) \sum_{\ialpha =1_{\alpha}}^{\lalpha}\sum_{\jalpha =1_{\alpha}}^{\lalpha} \big(G^{\alpha}\big)^{\ialpha \jalpha}\frac{\partial W}{\partial x^{\ialpha} }\frac{\partial W}{\partial x^{\jalpha}}=E=:a_{1}.
\end{gather}
Now, since
\[
\sum_{\alpha =1}^{r}\left(\frac{s^{\alpha 1}}{\det S}\right)s_{\alpha \beta}=\delta_{1 \beta},
\]
we may rewrite (\ref{HJPainform}) as
\begin{gather} \label{HJ4}
\sum_{\alpha =1}^{r}\frac{s^{\alpha 1}}{\det S}\left(\sum_{\ialpha =1_{\alpha}}^{\lalpha}\sum_{\jalpha =1_{\alpha}}^{\lalpha}\big(G^{\alpha}\big)^{\ialpha \jalpha}\frac{\partial W}{\partial x^{\ialpha} }\frac{\partial W}{\partial x^{\jalpha}}-\sum_{\beta=1}^{r}s_{\alpha \beta}a_{\beta}\right)=0,
\end{gather}
for any $(a_1,\dots,a_r) \in U \subset \mathbb{R}^r$. Choosing $W$ in the block-separable form~(\ref{Painsepform}), we see that any solution of the reduced Hamilton--Jacobi equations
\begin{gather} \label{HJ5}
\sum_{\ialpha =1_{\alpha}}^{\lalpha}\sum_{\jalpha =1_{\alpha}}^{\lalpha}\big(G^{\alpha}\big)^{\ialpha \jalpha}\frac{\partial W_{\alpha}}{\partial x^{\ialpha} }\frac{\partial W_{\alpha}}{\partial x^{\jalpha}}=\sum_{\beta = 1}^{r}s_{\alpha \beta}a_{\beta},
\end{gather}
will provide a solution of (\ref{HJ4}). But the latter equation always admit \emph{locally} solutions $W_\alpha(\textbf{x}^\alpha, \allowbreak a_1,\dots,a_r)$ by standard PDE results \cite{Tay2011a}. Differentiating (\ref{HJ5}) with respect to $a_{\beta}$, we obtain
\begin{gather*} %\label{HJ6}
2\sum_{\alpha =1}^{r}\sum_{\ialpha =1_{\alpha}}^{\lalpha}\sum_{\jalpha =1_{\alpha}}^{\lalpha}\big(G^{\alpha}\big)^{\ialpha \jalpha}\left(\frac{\partial W_{\alpha}}{\partial x^{\ialpha}}\right)\left( \frac{\partial^{2}W_{\alpha}}{\partial a_{\beta}\partial x^{\jalpha}}\right)=s_{\alpha \gamma}.
\end{gather*}
Since the generalized St\"ackel matrix (\ref{genStackelmat}) is to be non-singular, it follows that the rank condition that must be satisfied by our block-separable parametrized family of solutions of the Hamilton--Jacobi equation is precisely~(\ref{rankPainHJ}), thus proving that metrics of the Painlev\'e form~(\ref{Painleveform}) are indeed characterized by the existence of a parametrized family of solutions of the Hamilton--Jacobi equation satisfying a suitable completeness condition.
\end{proof}

Let us now give another proof of Proposition~\ref{HJSeparability-Painleve} that will allow us to characterize Painlev\'e metrics in terms of the generalized Levi-Civita conditions (\ref{genLeviCivita}). Working in block-orthogonal coordinates or directly with a Painlev\'e metric with the identification
\begin{gather}\label{not}
c^\alpha:=\frac{\det S}{s^{\alpha 1}},
\end{gather}
the Hamilton--Jacobi equation (\ref{HJ}) takes the form
\begin{gather}\label{HJPainformnew}
\sum_{\alpha =1}^{r} c^\alpha \sum_{\ialpha =1_{\alpha}}^{\lalpha}\sum_{\jalpha =1_{\alpha}}^{\lalpha} \big(G^{\alpha}\big)^{\ialpha \jalpha}\frac{\partial W}{\partial x^{\ialpha} }\frac{\partial W}{\partial x^{\jalpha}}=E=:a_{1}.
\end{gather}
We recall that we seek a solution $W$ of (\ref{HJPainformnew}) which is additively separable into groups of variables, that is
\[
W=\sum_{\alpha =1}^{r}W_{\alpha}\big({\bf x}^{\alpha}\big),
\]
and let
\[
u_{\alpha}^{(1)}=\sum_{\ialpha =1_{\alpha}}^{\lalpha}\sum_{\jalpha =1_{\alpha}}^{\lalpha}\big(G^{\alpha}\big)^{\ialpha \jalpha}\frac{\partial W_{\alpha}}{\partial x^{\ialpha} }\frac{\partial W_{\alpha}}{\partial x^{\jalpha}},
\]
in which case the Hamilton--Jacobi equation (\ref{HJPainformnew}) takes the form
\[
\sum_{\alpha =1}^{r}c^{\alpha}u_{\alpha}^{(1)}=a_{1}.
\]
Differentiating the latter equation with respect to $x^{\kbeta}$, we obtain the first order differential system in normal form
\begin{gather}\label{normalsyst}
 \partial_{\kbeta}u^{(1)}_{\gamma} =   \begin{cases}\displaystyle -\frac{1}{c^{\beta}}\sum_{\alpha =1}^{r}\big(\partial_{\kbeta} c^{\alpha}\big) u^{(1)}_{\alpha} & \gamma = \beta, \\
 0 & \gamma \neq \beta . \end{cases}
\end{gather}
Introducing the first-order differential operators
\begin{gather}\label{totaldiff}
D_{\kbeta} := \frac{\partial}{\partial x^{\kbeta}}-\frac{1}{c^{\beta}}\sum_{\alpha =1}^{r}\big(\partial_{\kbeta} c^{\alpha}\big) u^{(1)}_{\alpha}\frac{\partial}{\partial u^{(1)}_{\beta}},
\end{gather}
the differential system (\ref{normalsyst}) will admit a family of solutions $u^{(1)}_{\alpha}=u^{(1)}_{\alpha}\big(x^{1},\dots,x^{n};a_{1},\dots,a_{r}\big)$ defined on an open subset $U\in M$ and depending smoothly on~$r$ constants $(a_{1},\dots,a_{r})$ defined in an open subset $A\subset {\mathbb{R}}^{r}$, satisfying the rank condition
\begin{gather}\label{rankHJ}
\det\left(\frac{\partial u^{(1)}_{\alpha}}{\partial a_{\beta}}\right)\neq 0,
\end{gather}
if and only if the operators (\ref{totaldiff}) pairwise commute, that is
\begin{gather}\label{commutation}
\big[D_{\kbeta},D_{\jalpha}\big]=0
\end{gather}
for all $1\leq \alpha, \beta \leq r$, $1_{\alpha}\leq \jalpha\leq \lalpha$, $1_{\beta}\leq \kbeta\leq \lbeta$. We refer to \cite[Theorem~2.1]{BCR1-2002} for this result. Note that if~(\ref{commutation}) hold, then the Hamilton--Jacobi equation admits locally a solution which is additively separable into groups of variables and satisfies the completeness condition~(\ref{rankPainHJ}) as a consequence of the completeness condition (\ref{rankHJ}). In consequence, such metrics are of the Painlev\'e form~(\ref{Painleveform}).

We now prove that the commutation relations (\ref{commutation}) are equivalent to the generalized Levi-Civita conditions~(\ref{genLeviCivita}). Note that this is a natural generalization of the link between the complete separation of variables which is familiar from the St\"ackel case and the classical Levi-Civita separability conditions, as reviewed in~\cite{BCR1-2002,KaMi1984}. Indeed, we have
\begin{Lemma}
The pairwise commutation relations \eqref{commutation} for the derivations $D_{\kbeta}$ are equivalent to the generalized Levi-Civita conditions~\eqref{genLeviCivita}.
\end{Lemma}

\begin{proof}Substituting the expression (\ref{totaldiff}) of the differential operators $D_{\kbeta}$ into the commutation relations (\ref{commutation}), we obtain
 \begin{gather*}
 \left[-\partial_{\kbeta}\left(\frac{1}{c^{\gamma}}\sum_{\delta =1}^{r}\big(\partial_{\pgamma} c^{\delta}\big)u^{(1)}_{\delta}\right) + \frac{1}{c^{\beta} c^{\gamma}}\sum_{\alpha =1}^{r}\big(\partial_{\kbeta} c^{\alpha}\big)u^{(1)}_{\alpha}\big(\partial_{\pgamma} c^{\beta}\big)\right]\frac{\partial}{\partial u^{(1)}_{\gamma}} \\
 \qquad{} +
 \left[-\partial_{\kbeta}\left(\frac{1}{c^{\gamma}}\sum_{\delta =1}^{r}\big(\partial_{\pgamma} c^{\delta}\big)u^{(1)}_{\delta}\right)+\frac{1}{c^{\beta}c^{\gamma}}\sum_{\alpha =1}^{r}\big(\partial_{\kbeta}c^{\alpha}\big)u^{(1)}_{\alpha}\big(\partial_{\pgamma} c^{\beta}\big)\right]\frac{\partial}{\partial u^{(1)}_{\gamma}}=0.
 \end{gather*}
For $1\leq \gamma \neq \beta \leq r$, the above identity is equivalent to
\begin{gather} \label{GenLC}
\big(\partial_{\kbeta} \log c^{\gamma}\big)\big(\partial_{\pgamma} \log c^{\alpha}\big)+\big(\partial_{\kbeta} \log c^{\alpha}\big)\big(\partial_{\pgamma} \log c^{\beta}\big)-\frac{1}{c^{\alpha}}\partial_{\kbeta}\partial_{\pgamma}c^{\alpha}=0,
\end{gather}
which are precisely the generalized Levi-Civita conditions~(\ref{genLeviCivita}) after relabeling of indices. Note that for $1\leq \gamma = \beta \leq r$, the above identity is always satisfied by a straightforward calculation.
\end{proof}

At this stage, we thus have proved another characterization of Painlev\'e metrics which appears in \cite[p.~12]{CR2019}.

\begin{Proposition} \label{PainLC}A metric $g$ is of Painlev\'e type if and only if there exist block orthogonal coordinates~$\big(\textbf{x}^\alpha\big)$ such that the generalized Levi-Civita conditions \eqref{GenLC} hold.
\end{Proposition}

We finish this section giving still another characterization of Painlev\'e metrics of a more intrinsic nature. The starting point is the observation that the generalized Killing--Eisenhart equations are related to the existence of quadratic first integrals (or symmetries) $K_{(\beta)}$ by the following result proved in~\cite[Proposition~5.3]{CR2019} (see also Lemma~\ref{KE} for an implicit proof of this proposition)

\begin{Proposition}In block orthogonal coordinates, we have that
\[
 \big\{H,K_{(\beta)} \big\} = 0,
\]	
if and only if the Killing--Eisenhart equations
\begin{gather}\label{genKE1}
 \partial_{j_{\gamma}}(\rho_{ \beta \delta})= (\rho_{\beta \gamma} - \rho_{\beta \delta }) \big(\partial_{j_{\gamma}} \log c^\delta\big),
\end{gather}
hold for all $1 \leq \gamma, \delta \leq r$.
\end{Proposition}

The second observation is the fact that the integrability conditions for the Killing--Eisenhart equations~(\ref{genKE1}) are given by
\begin{gather*} %\label{IC1}
 (\rho_{\beta \gamma} - \rho_{\beta \delta}) \left[ \big( \partial_{\jgamma} \log c^{\delta} \big) \big( \partial_{\kalpha} \log c^{\gamma}\big) + \big(\partial_{\jgamma} \log c^{\alpha} \big) \big( \partial_{\kalpha} \log c^{\delta} \big) - \frac{1}{c^{\delta}} \partial_{\kalpha}\partial_{\jgamma}c^{\delta} \right] = 0,
\end{gather*}
for all $1 \leq \alpha, \beta, \gamma, \delta \leq r$. Clearly, these integrability conditions can be shortened as
\begin{gather} \label{IC2}
 (\rho_{\beta \gamma} - \rho_{\beta \delta}) \cdot \left[ \textrm{Levi-Civita conditions} \right] = 0.
\end{gather}

Using these two observations, Chanu ad Rastelli proved in \cite[Proposition 5.5]{CR2019} the following characterization of Painlev\'e metrics.

\begin{Proposition} \label{PainKA}
$(M,g)$ is a Painlev\'e manifold if and only if
\begin{enumerate}\itemsep=0pt
\item[$1)$] there exist $r$ independent quadratic first integral $K_{(\beta)}$, $\beta = 1,\dots,r$ such that $K_{(1)} = H$ and $\{ H, K_{(\beta)} \} = 0$.
\item[$2)$] The associated Killing two tensors	$K_{(\beta)}^{ij}$ are simultaneously block-diagonalized and have common normally integrable eigenspaces.
\end{enumerate}
\end{Proposition}
\begin{proof}
If $g$ is a Painlev\'e metric, then the above assertions were already proved.

Assume now that there exist $r$ linearly independent Killing tensors simultaneously in block diagonal form. Then the generalized Killing--Eisenhart equations (\ref{genKE1}) will admit an $r$-dimen\-sio\-nal vector space of solutions. The latter is equivalent to the invertibility of the fundamental matrix $\mathcal{A}$ defined by
\[
{\mathcal{A}}=\left(
\begin{matrix}
\rho_{11}&\dots&\rho_{1r} \\
\vdots& & \vdots \\
\rho_{r1}&\dots&\rho_{rr}
\end{matrix}
\right).
\]
Hence any solution $(\rho_{1},\dots,\rho_{r})$ of (\ref{genKE1}) is given by
\[
\left(
\begin{matrix}\rho_{1}\\
\vdots\\
\rho_{r}
\end{matrix}
\right)
=
{\mathcal{A}}\left(
\begin{matrix}a_{1}\\
\vdots\\
a_{r}
\end{matrix}
\right),
\]
for some constants $(a_{1},\dots,a_{r})\in A$. It is clear then that the Killing--Eisenhart equations are completely integrable and thus satisfy the integrability conditions (\ref{IC2}). Moreover, from the invertibility of~$\mathcal{A}$, we can always choose the constants $(a_1,\dots,a_r)$ such that $\rho_{\alpha} \neq \rho_{\beta}$, $\forall\, 1 \leq \alpha \ne \beta \leq r$ at a point $p\in M$ and therefore in an neighbourhood of $p$, by continuity. Hence, the integrabi\-li\-ty conditions~(\ref{IC2}) reduces to the generalized Levi-Civita separability conditions~(\ref{GenLC}). We conclude that $g$ is a Painlev\'e metric from Proposition~\ref{PainLC}.
\end{proof}

As a concluding remark for this section, we emphasize that the hypotheses~(\ref{H}) that we make to deduce the Levi-Civita conditions from the Killing--Eisenhart equations aren't in fact necessary. Indeed, it follows from Proposition~\ref{PainLC} or Proposition~\ref{PainKA} that the Levi-Civita conditions always hold whenever~$g$ is a Painlev\'e metric, that is a metric of the form~(\ref{Painleveform}).


\setcounter{section}{5}
\setcounter{subsection}{0}
\setcounter{subsubsection}{0}
\setcounter{Theorem}{2}
\setcounter{equation}{8}
% Original source lines 1121--1123
\section{Proofs of the main theorems} \label{ProofsofMainTheorems}

\subsection{Proof of Theorem \ref{maintheorem1}}

\setcounter{section}{6}
\setcounter{subsection}{1}
\setcounter{subsubsection}{0}
\setcounter{Theorem}{0}
\setcounter{equation}{0}
% Original source lines 1126--1318
What remains to be done in order to prove Theorem \ref{maintheorem1} and what constitutes our main task is therefore to show that the generalized Robertson conditions are equivalent to the conditions~(\ref{Ricciform}) on the Ricci tensor, thus generalizing the classical result of Eisenhart \cite{Eis1934} to Painlev\'e metrics. In order to do this, we will show that
\begin{gather}\label{adaptedRicci}
R_{\jalpha \kbeta}=-\frac{3}{4}\partial_{j_{\alpha}}\partial_{k_{\beta}}\log\left[\frac{(\det S)^{n-2}}{\big(s^{11}\big)^{l_{1}}\cdots \big(s^{r1}\big)^{l_{r}}}\right] + \frac{1}{4}T_{\jalpha \kbeta},
\end{gather}
where
\begin{gather}
T_{\jalpha \kbeta} =(l_{\alpha}+l_{\beta}-2)\frac{s^{\alpha 1}s^{\beta 1}}{\det S}\partial_{j_{\alpha}}\partial_{k_{\beta}}\left(\frac{\det S}{s^{\alpha 1}s^{\beta 1}}\right)\nonumber\\
\hphantom{T_{\jalpha \kbeta} =}{} +\sum_{\gamma\neq \alpha,\beta=1}^{r}l_{\gamma}\left[\left(\partial_{j_{\alpha}}\log\frac{s^{\gamma 1}}{\det S}\right)\left(\partial_{k_{\beta}}\log\frac{s^{\alpha 1}}{\det S}\right)\right.\nonumber\\
\left.\hphantom{T_{\jalpha \kbeta} =}{}+\left(\partial_{j_{\alpha}}\log\frac{s^{\beta 1}}{\det S}\right)\left(\partial_{k_{\beta}}\log\frac{s^{\gamma 1}}{\det S}\right)
-\frac{\det S}{s^{\gamma 1}}\partial_{j_{\alpha}}\partial_{k_{\beta}}\left(\frac{s^{\gamma 1}}{\det S}\right)\right].\label{Tterms}
\end{gather}
We note that the expression (\ref{adaptedRicci}) of the off block-diagonal components of the Ricci tensor $R_{\jalpha \kbeta}$ is independent of the $r$ Riemannian block metrics $G_{\beta}$, $1\leq \beta \leq r$ defined by~(\ref{blockmetric}). We also remark that the first term in the expression~(\ref{Tterms}) of~$T_{\jalpha \kbeta}$, involving second derivatives, vanishes identically in the special case of St\"ackel metrics since the pre-factor $\lalpha+\lbeta-2$ is zero in that case.

Once we will have established (\ref{adaptedRicci}), it will then follow from Lemma~\ref{LeviCivita} and more precisely from the generalized Levi-Civita conditions~(\ref{genLeviCivita}) and~(\ref{secondorder}) that the generalized Robertson conditions (\ref{GenRobertson}) are indeed equivalent to the vanishing conditions~(\ref{Ricciform}) on the non-block diagonal components of the Ricci tensor. We therefore proceed to establish the form (\ref{adaptedRicci}) of the Ricci tensor for a Painlev\'e metric.

The expression of the Ricci tensor in terms of the Christoffel symbols is given by
\begin{gather}\label{Ricci}
R_{\jalpha \kbeta}=R^{l}_{{\phantom l} l \jalpha \kbeta}=\partial_{l}{\Gamma^{l}}_{\jalpha \kbeta}-\partial_{\jalpha}{\Gamma^{l}}_{l\kbeta}+{\Gamma^{m}}_{\jalpha \kbeta}{\Gamma^{l}}_{l m}-{\Gamma^{m}}_{l\kbeta}{\Gamma^{l}}_{\jalpha m},
\end{gather}
where the summation convention is applied with $1\leq l,m \leq n = \dim M$. In order to compute the right-hand side of (\ref{Ricci}), we will need expressions for the Christoffel symbols of a Painlev\'e metric (\ref{Painleveform}). Using the standard formulas
\begin{gather*}%\label{Christoffelgen}
\Gamma_{hji}=\frac{1}{2}(\partial_{h}g_{ji}+\partial_{j}g_{ih}-\partial_{i}g_{hj}), \qquad {\Gamma^{i}}_{hj}=g^{ik}\Gamma_{hjk},
\end{gather*}
and writing the Painlev\'e metric (\ref{Painleveform}) in block-diagonal form as
\[
{\rm d}s^{2}=\sum_{\alpha = 1}^{r}\sum_{\ialpha=1_{\alpha}}^{\lalpha}\sum_{\jalpha=1_{\alpha}}^{\lalpha}(g_{\alpha})_{\ialpha \jalpha}{\rm d}x^{\ialpha}{\rm d}x^{\jalpha},
\]
we obtain for fixed indices $1\leq \alpha,\beta \leq r$,
\begin{gather}
{\Gamma^{\ialpha}}_{\kalpha \jbeta}=\frac{1}{2}\sum_{\palpha=1_{\alpha}}^{\lalpha} \big(g^{\alpha}\big)^{\ialpha \palpha}\partial_{\jbeta}(g_{\alpha})_{\kalpha \palpha},\nonumber\\
 {\Gamma^{\jbeta}}_{\ialpha \kalpha}=-\frac{1}{2}\sum_{\pbeta=1_{\beta}}^{\lbeta}\big(g^{\beta}\big)^{\jbeta \pbeta}\partial_{\pbeta}(g_{\alpha})_{\ialpha \kalpha}\qquad \text{for} \quad \alpha \neq \beta,\label{Christoffel}
\end{gather}
and
\begin{gather}\label{Christoffelmixed}
{\Gamma^{\ialpha}}_{\halpha \kalpha}=\frac{1}{2}\sum_{\kalpha=1_{\alpha}}^{\lalpha} g^{\ialpha \kalpha}(\partial_{\halpha}g_{\jalpha \kalpha}+\partial_{\jalpha}g_{\halpha \kalpha}-\partial_{\kalpha}g_{\halpha \jalpha}).
\end{gather}
In view of the expressions (\ref{Christoffel}) of the Christoffel symbols, it is convenient to split the sum over~$l$ appearing in (\ref{Ricci}) into three sums, the first sum corresponding to the values of the summation index $l$ lying in groups of indices different from the groups corresponding to $\alpha$ and $\beta$, and the remaining two to the values of $l$ belonging to the groups of indices labelled by $\alpha$ and $\beta$ respectively. Thus we write
\begin{gather}\label{Riccidecomp}
R_{\jalpha \kbeta}=\sum_{l=1}^{n}R^{l}_{{\phantom l} l \jalpha \kbeta}=\sum_{\gamma \neq \alpha,\, \beta=1}^{r}\sum_{\pgamma=1_{\gamma}}^{\lgamma}R^{\pgamma}_{{\phantom l} \pgamma \jalpha \kbeta} + \sum_{\palpha =1_{\alpha}}^{\lalpha}R^{\palpha}_{{\phantom l} \palpha \jalpha \kbeta}+\sum_{\pbeta=1_{\beta}}^{\lbeta}R^{\pbeta}_{{\phantom l} \pbeta \jalpha \kbeta}.
\end{gather}
Let us begin with the first term. We have, for $\gamma\neq \alpha, \beta$,
\begin{gather}
\sum_{\pgamma = 1_{\gamma}}^{\lgamma}R^{\pgamma}_{{\phantom l} \pgamma \jalpha \kbeta}=-\sum_{\pgamma = 1_{\gamma}}^{\lgamma}\partial_{\jalpha}{\Gamma^{\pgamma}}_{\pgamma \kbeta} +\sum_{\pgamma = 1_{\gamma}}^{\lgamma}\sum_{\malpha=1_{\alpha}}^{\lalpha}{\Gamma^{\malpha}}_{\jalpha \kbeta}{\Gamma^{\pgamma}}_{\pgamma \malpha}\nonumber\\
\hphantom{\sum_{\pgamma = 1_{\gamma}}^{\lgamma}R^{\pgamma}_{{\phantom l} \pgamma \jalpha \kbeta}=}{}
+\sum_{\pgamma = 1_{\gamma}}^{\lgamma}\sum_{\mbeta=1_{\beta}}^{\lbeta}{\Gamma^{\mbeta}}_{\jalpha \kbeta}{\Gamma^{\pgamma}}_{\pgamma \mbeta} -\sum_{\pgamma = 1_{\gamma}}^{\lgamma}\sum_{\mgamma=1_{\gamma}}^{\lgamma}{\Gamma^{\mgamma}}_{\pgamma \kbeta}{\Gamma^{\pgamma}}_{\jalpha \mgamma},\label{Riemann1}
\end{gather}
where the summations have been written out explicitly to avoid notational ambiguities. We begin with evaluating the first-derivative term on the right-hand side of~(\ref{Riemann1}) using the expressions~(\ref{Christoffel}) for the Christoffel symbols, thus obtaining, for $\gamma \neq \alpha, \beta$,
\begin{gather}\label{derivgamma}
\sum_{\pgamma = 1_{\gamma}}^{\lgamma}\partial_{\jalpha}{\Gamma^{\pgamma}}_{\pgamma \kbeta} = \frac{1}{2}\sum_{\pgamma = 1_{\gamma}}^{\lgamma}\sum_{\ngamma = 1_{\gamma}}^{\lgamma}\partial_{\jalpha}\big(\big(g^{\gamma}\big)^{\pgamma \ngamma}\partial_{\kbeta}((g_{\gamma})_{\pgamma \ngamma})\big)=\frac{1}{2}\partial_{\jalpha}\partial_{\kbeta}\log(|g_{\gamma}|),
\end{gather}
where $|g_{\gamma}|:=\det(g_{\igamma \jgamma})$. It follows that
\[
\sum_{\gamma \neq \alpha,\, \beta=1}^{r}\sum_{\pgamma=1_{\gamma}}^{\lgamma}\partial_{\jalpha}{\Gamma^{\pgamma}}_{\pgamma \kbeta}= \frac{1}{2}\partial_{\jalpha}\partial_{\kbeta}\left(\prod_{\gamma \neq \alpha, \,\beta=1}^{r}\log(|g_{\gamma}|)\right).
\]
From the Painlev\'e form (\ref{Painleveform}), we have
\begin{gather}\label{Painleveblock}
g_{\igamma \jgamma}=\frac{\det S}{s^{\gamma 1}}(G_{\gamma})_{\igamma \jgamma},
\end{gather}
so that
\[
|g_{\gamma}|=\frac{(\det S)^{\lgamma}}{\big({s^{\gamma 1}}\big)^{\lgamma}}|G_{\gamma}|.
\]
Using the fact that $|G_{\gamma}|$ is a function of the variables ${\bf x}^{\gamma}$ only, we obtain
\begin{gather*}%\label{secondderiv1}
\sum_{\gamma \neq \alpha, \,\beta=1}^{r}\sum_{\pgamma=1_{\gamma}}^{\lgamma} \partial_{\jalpha}{\Gamma^{\pgamma}}_{\pgamma \kbeta}=\frac{1}{2}\sum_{\gamma \neq \alpha, \,\beta=1}^{r}\lgamma \partial_{\jalpha}\partial_{\kbeta}\left(\log\left(\frac{\det S}{s^{\gamma 1}}\right)\right).
\end{gather*}
We notice that the above expression is independent of the quadratic differential forms $G_{\alpha}$ defined by~(\ref{blockmetric}).

Next we evaluate the terms quadratic in the Christoffel symbols in the right-hand side of~(\ref{Riemann1}). Again, we use the fact that $\gamma\neq \alpha, \beta$ and the fact that in the Painlev\'e form~(\ref{Painleveform}), each of the quadratic differential forms $G_{\alpha}$ defined by~(\ref{blockmetric}) depends on the group of variables ${\bf x}^{\alpha}$ only. We have
\begin{gather}
\sum_{\malpha=1_{\alpha}}^{\lalpha} \sum_{\pgamma = 1_{\gamma}}^{\lgamma}{\Gamma^{\malpha}}_{\jalpha \kbeta}{\Gamma^{\pgamma}}_{\pgamma \malpha}\nonumber\\
\qquad{} =\frac{1}{4}\sum_{\malpha = 1_{\alpha}}^{l_{\alpha}} \sum_{\pgamma = 1_{\gamma}}^{\lgamma}\sum_{\nalpha=1_{\alpha}}^{l_{\alpha}}(g^{\alpha})^{\malpha \nalpha}\partial_{\kbeta}((g_{\alpha})_{\jalpha \nalpha})g^{\pgamma \hgamma}\partial_{\malpha}((g_{\gamma})_{\pgamma \hgamma})\nonumber\\
\qquad{} =\frac{1}{4}\sum_{\malpha=1_{\alpha}}^{l_{\alpha}}\sum_{\nalpha=1_{\alpha}}^{l_{\alpha}}(g^{\alpha})^{\malpha \nalpha}\partial_{\kbeta}((g_{\alpha})_{\jalpha \nalpha})\partial_{\malpha}\log(|g_{\gamma}|)\nonumber\\
\qquad{} =\frac{1}{4}\sum_{\malpha=1_{\alpha}}^{l_{\alpha}}\sum_{\nalpha=1_{\alpha}}^{l_{\alpha}}\frac{s^{\alpha 1}}{\det S}\big(G^{\alpha}\big)^{\malpha \nalpha}\partial_{\kbeta}\left(\frac{\det S}{s^{\alpha 1}}(G_{\alpha})_{\jalpha \nalpha}\right)\partial_{\malpha}\log(|g_{\gamma}|)\nonumber\\
\qquad{} =\frac{1}{4}\sum_{\malpha=1_{\alpha}}^{l_{\alpha}}\sum_{\nalpha=1_{\alpha}}^{l_{\alpha}}{\delta^{\malpha}}_{\jalpha}\partial_{\kbeta}\log\left(\frac{\det S}{s^{\alpha 1}}\right)\partial_{\malpha}\log \left(\frac{(\det S)^{\lgamma}}{(s^{\gamma 1})^{\lgamma}}\right)\nonumber\\
\qquad{} =\frac{1}{4}\left(\partial_{\kbeta}\log\left(\frac{s^{\alpha 1}}{\det S}\right)\right)\left(\partial_{\jalpha}\log \left(\frac{\big(s^{\gamma 1}\big)^{\lgamma}}{(\det S)^{\lgamma}}\right)\right).\label{quadr1}
\end{gather}
We notice that the above expression is again independent of the quadratic differential forms $G_{\alpha}$ defined by (\ref{blockmetric}). Likewise, we obtain for the next quadratic term in the Christoffel symbols that appears in the right-hand side of~(\ref{Riemann1}),
\begin{gather}\label{quadr2}
\sum_{\mbeta=1_{\beta}}^{\lbeta} \sum_{\pgamma = 1_{\gamma}}^{\lgamma}{\Gamma^{\mbeta}}_{\jalpha \kbeta}{\Gamma^{\pgamma}}_{\pgamma \mbeta}=\frac{1}{4}\left(\partial_{\jalpha}\log\left(\frac{s^{\beta 1}}{\det S}\right)\right)\left(\partial_{\kbeta}\log \left(\frac{\big(s^{\gamma 1}\big)^{\pgamma}}{(\det S)^{\pgamma}}\right)\right).
\end{gather}
For the third and final quadratic term, we have
\begin{align}
\sum_{\mgamma=1_{\gamma}}^{\lgamma}{\Gamma^{\mgamma}}_{\pgamma \kbeta}{\Gamma^{\pgamma}}_{\jalpha \mgamma} &=\frac{1}{4}\sum_{\mgamma=1_{\gamma}}^{\lgamma}{\delta^{\mgamma}}_{\pgamma}\left(\partial_{\kbeta}\log\left(\frac{s^{\gamma 1}}{\det S}\right)\right){\delta^{\pgamma}}_{\mgamma}\left(\partial_{\jalpha}\log\left(\frac{s^{\gamma 1}}{\det S}\right)\right)\nonumber\\
&=\frac{1}{4}\lgamma \left(\partial_{\kbeta}\log\left(\frac{s^{\gamma 1}}{\det S}\right)\right)\left(\partial_{\jalpha}\log\left(\frac{s^{\gamma 1}}{\det S}\right)\right),\label{quadr3}
\end{align}
which is likewise independent of the quadratic differential forms $G_{\alpha}$ defined by (\ref{blockmetric}). Putting together the expressions (\ref{derivgamma}), (\ref{quadr1}), (\ref{quadr2}) and (\ref{quadr3}) we obtain
\begin{gather}
\sum_{1\leq \gamma \neq \alpha, \beta\leq r}R^{\pgamma}_{{\phantom l} \pgamma \jalpha \kbeta}=\sum_{\gamma \neq \alpha, \,\beta =1}^{r}l_{\gamma}\left[\frac{1}{2}\partial_{\jalpha}\partial_{\kbeta}\log\left(\frac{s^{\gamma 1}}{\det S}\right)\right.\nonumber\\
\hphantom{\sum_{1\leq \gamma \neq \alpha, \beta\leq r}R^{\pgamma}_{{\phantom l} \pgamma \jalpha \kbeta}=}{}
+\frac{1}{4}\left(\partial_{\kbeta}\log\left(\frac{s^{\alpha 1}}{\det S}\right)\right)\left(\partial_{\jalpha}\log\left(\frac{s^{\gamma 1}}{\det S}\right)\right)\nonumber\\
\hphantom{\sum_{1\leq \gamma \neq \alpha, \beta\leq r}R^{\pgamma}_{{\phantom l} \pgamma \jalpha \kbeta}=}{}
+ \frac{1}{4}\left(\partial_{\jalpha}\log\left(\frac{s^{\beta 1}}{\det S}\right)\right)\left(\partial_{\kbeta}\log\left(\frac{s^{\gamma 1}}{\det S}\right)\right)\nonumber\\
\left.\hphantom{\sum_{1\leq \gamma \neq \alpha, \beta\leq r}R^{\pgamma}_{{\phantom l} \pgamma \jalpha \kbeta}=}{}
-\frac{1}{4}\left(\partial_{\kbeta}\log\left(\frac{s^{\gamma 1}}{\det S}\right)\right)\left(\partial_{\jalpha}\log\left(\frac{s^{\gamma 1}}{\det S}\right)\right)\right].\label{Riccidecomp1}
\end{gather}
We still need to evaluate the curvature components $R^{\palpha}_{{\phantom l} \palpha \jalpha \kbeta}$ and $R^{\pbeta}_{{\phantom l} \pbeta \jalpha \kbeta}$, which will require a~separate calculation. We have
\begin{gather}
\sum_{\palpha=1_{\alpha}}^{\lalpha}R^{\palpha}_{{\phantom l} \palpha \jalpha \kbeta} =\sum_{\palpha=1_{\alpha}}^{\lalpha}\partial_{\palpha}{\Gamma^{\palpha}}_{\jalpha \kbeta}-\sum_{\palpha=1_{\alpha}}^{\lalpha}\partial_{\jalpha}{\Gamma^{\palpha}}_{\palpha \kbeta}\nonumber\\
\hphantom{\sum_{\palpha=1_{\alpha}}^{\lalpha}R^{\palpha}_{{\phantom l} \palpha \jalpha \kbeta} =}{}
+\sum_{\palpha=1_{\alpha}}^{\lalpha}\sum_{\malpha = 1_{\alpha}}^{\lalpha}{\Gamma^{\malpha}}_{\jalpha \kbeta}{\Gamma^{\palpha}}_{\palpha \malpha}+\sum_{\palpha=1_{\alpha}}^{\lalpha}\sum_{\mbeta=1_{\beta}}^{\lbeta}{\Gamma^{\mbeta}}_{\jalpha \kbeta}{\Gamma^{\palpha}}_{\palpha \mbeta}\nonumber\\
\hphantom{\sum_{\palpha=1_{\alpha}}^{\lalpha}R^{\palpha}_{{\phantom l} \palpha \jalpha \kbeta} =}{}-\sum_{\palpha=1_{\alpha}}^{\lalpha}\sum_{\malpha=1_{\alpha}}^{\lalpha}{\Gamma^{\malpha}}_{\palpha \kbeta}{\Gamma^{\palpha}}_{\jalpha \malpha}-\sum_{\palpha=1_{\alpha}}^{\lalpha}\sum_{\mbeta=1_{\beta}}^{\lbeta}{\Gamma^{\mbeta}}_{\lalpha \kbeta}{\Gamma^{\palpha}}_{\jalpha \mbeta},\label{Riemann2}
\end{gather}
where again we have written out the summation signs explicitly to avoid notational ambiguities. We have, using (\ref{Christoffel}),
\begin{gather}\label{moreChristoffel}
{\Gamma^{\palpha}}_{\jalpha \kbeta}=\frac{1}{2}{\delta^{\palpha}}_{\jalpha}\partial_{\kbeta}\left(\log\left(\frac{\det S}{s^{\alpha 1}}\right)\right),
\end{gather}
so using the fact that the cofactor $s^{\alpha 1}$ is independent of ${\bf x}^{\alpha}$, we obtain
\begin{gather}\label{newderiv1}
\sum_{\palpha=1_{\alpha}}^{\lalpha}\partial_{\palpha}{\Gamma^{\palpha}}_{\jalpha \kbeta}=\frac{1}{2}\partial_{\jalpha}\partial_{\kbeta}\big(\log(\det S)\big),
\end{gather}
and a similar calculation gives
\begin{gather}\label{newderiv2}
\sum_{\palpha=1_{\alpha}}^{\lalpha}\partial_{\jalpha}{\Gamma^{\palpha}}_{\palpha \kbeta}=\frac{1}{2}\partial_{\jalpha}\partial_{\kbeta}\big(\log\big((\det S)^{\lalpha}\big)\big).
\end{gather}
We now evaluate the quadratic terms in the Christoffel symbols that appear in the curvature component~(\ref{Riemann2}). In order to do so, we substitute into the expression (\ref{Christoffelmixed}) of the Christoffel symbols the expressions
\[
(g_{\alpha})_{\ialpha \jalpha}=\frac{\det S}{s^{\alpha 1}}G_{\alpha},
\]
which result from the Painlev\'e form (\ref{Painleveform}). We obtain the following expressions for the Christoffel symbols,
\begin{gather}
{\Gamma^{\ialpha}}_{\halpha \jalpha}={\gamma^{\ialpha}}_{\halpha \jalpha} +\frac{1}{2}\Bigg[ {\delta^{\ialpha}}_{\jalpha}\partial_{\halpha}\big(\log{(\det S)}\big)+{\delta^{\ialpha}}_{\halpha}\partial_{\jalpha}\big(\log{(\det S)}\big)\nonumber\\
\hphantom{{\Gamma^{\ialpha}}_{\halpha \jalpha}=}{}-\sum_{\kalpha = 1_{\alpha}}^{\lalpha}\big(G^{\alpha}\big)^{\ialpha \kalpha}(G_{\alpha})_{\halpha \jalpha}\partial_{\kalpha}\big(\log{(\det S)}\big)\Bigg],\label{longChristoffel}
\end{gather}
where the ${\gamma^{\ialpha}}_{\halpha \jalpha}$ denote the Christoffel symbols of the block metric~$G_{\alpha}$ given by~(\ref{blockmetric}). It then follows from~(\ref{longChristoffel}) and~(\ref{moreChristoffel}) that
\begin{gather}
\sum_{\palpha=1_{\alpha}}^{\lalpha}\sum_{\malpha = 1_{\alpha}}^{\lalpha}{\Gamma^{\malpha}}_{\jalpha \kbeta}{\Gamma^{\palpha}}_{\palpha \malpha}\nonumber\\
\qquad{} =\frac{1}{2}\partial_{\kbeta}\left(\log\left(\frac{\det S}{s^{\alpha 1}}\right)\right)\left[\sum_{\palpha=1_{\alpha}}^{\lalpha}{\gamma^{\palpha}}_{\palpha \jalpha}+\frac{1}{2}\lalpha \partial_{\malpha}\big(\log{(\det S)}\big)\right].\label{newquadr1}
\end{gather}
Similarly, using (\ref{moreChristoffel}) and (\ref{longChristoffel}), we obtain
\begin{gather}
\sum_{\palpha=1_{\alpha}}^{\lalpha}\sum_{\malpha=1_{\alpha}}^{\lalpha}{\Gamma^{\malpha}}_{\palpha \kbeta}{\Gamma^{\palpha}}_{\jalpha \malpha}\nonumber\\
\qquad{} =\frac{1}{2}\partial_{\kbeta}\left(\log\left(\frac{\det S}{s^{\alpha 1}}\right)\right)\left[\sum_{\palpha=1_{\alpha}}^{\lalpha}{\gamma^{\palpha}}_{\jalpha \palpha}+\frac{1}{2}\lalpha \partial_{\malpha}\big(\log{(\det S)}\big)\right].\label{newquadr2}
\end{gather}
It follows therefore from (\ref{newquadr1}) and (\ref{newquadr2}) that
\[
\sum_{\palpha=1_{\alpha}}^{\lalpha}\sum_{\malpha = 1_{\alpha}}^{\lalpha}{\Gamma^{\malpha}}_{\jalpha \kbeta}{\Gamma^{\palpha}}_{\palpha \malpha}-\sum_{\palpha=1_{\alpha}}^{\lalpha}\sum_{\malpha=1_{\alpha}}^{\lalpha}{\Gamma^{\malpha}}_{\palpha \kbeta}{\Gamma^{\palpha}}_{\jalpha \malpha}=0.
\]
We now evaluate the remaining difference of two double sums in the expression (\ref{Riemann2}) of the curvature, using the expressions
\[
{\Gamma^{\mbeta}}_{\jalpha \kbeta}=\frac{1}{2}{\delta^{\mbeta}}_{\kbeta}\partial_{\jalpha}\left(\log\left(\frac{\det S}{s^{\beta 1}}  \right)\right), \qquad {\Gamma^{\mbeta}}_{\jalpha \kbeta} = \frac{1}{2}\partial_{\mbeta}\left(\log\left(\frac{(\det S)^{\lalpha}}{\big(s^{\beta 1}\big)^{\lalpha}} \right)\right),
\]
for the Christoffel symbols, to obtain
\begin{gather}
\sum_{\palpha=1_{\alpha}}^{\lalpha}\sum_{\malpha = 1_{\alpha}}^{\lalpha}{\Gamma^{\malpha}}_{\jalpha \kbeta}{\Gamma^{\palpha}}_{\palpha \malpha}-\sum_{\palpha=1_{\alpha}}^{\lalpha}\sum_{\mbeta=1_{\beta}}^{\lbeta}{\Gamma^{\mbeta}}_{\lalpha \kbeta}{\Gamma^{\palpha}}_{\jalpha \mbeta}\nonumber\\
\qquad{} =\frac{1}{4}(\lalpha-1)\partial_{\jalpha}\left(\log \left(\frac{\det S}{s^{\beta 1}}\right)\right)\partial_{\kbeta}\left(\log\left(\frac{\det S}{s^{\alpha 1}} \right)\right).\label{newquadr3}
\end{gather}
Substituting the expressions (\ref{newderiv1}), (\ref{newderiv2}), (\ref{newquadr1}), (\ref{newquadr2}) and (\ref{newquadr3}) into (\ref{Riemann2}), we obtain
\begin{gather}
\sum_{\palpha=1_{\alpha}}^{\lalpha}R^{\palpha}_{{\phantom l} \palpha \jalpha \kbeta}= \frac{1-\lalpha}{2} \partial_{\jalpha}\partial_{\kbeta} \big(\log(\det S)\big)\nonumber\\
\hphantom{\sum_{\palpha=1_{\alpha}}^{\lalpha}R^{\palpha}_{{\phantom l} \palpha \jalpha \kbeta}=}{}
+\frac{1}{4}(\lalpha-1)\partial_{\jalpha}\left(\log\left(\frac{\det S}{s^{\beta 1}}\right)\right)\partial_{\kbeta}\left(\log\left(\frac{\det S}{s^{\alpha 1}}\right)\right),\label{Riccidecomp2}
\end{gather}
and similarly
\begin{gather}\label{Riccidecomp3}
\sum_{\pbeta=1_{\beta}}^{\lalpha}R^{\pbeta}_{{\phantom l} \pbeta \jalpha \kbeta}=\frac{1-\lbeta}{2}\partial_{\jalpha}\partial_{\kbeta} \big(\log(\det S)\big)\nonumber\\
\hphantom{\sum_{\pbeta=1_{\beta}}^{\lalpha}R^{\pbeta}_{{\phantom l} \pbeta \jalpha \kbeta}=}{}
+\frac{1}{4}(\lbeta-1)\partial_{\jalpha}\left(\log\left(\frac{\det S}{s^{\beta 1}}\right)\right)\partial_{\kbeta}\left(\log\left(\frac{\det S}{s^{\alpha 1}}\right)\right).
\end{gather}
We now substitute the expressions (\ref{Riccidecomp1}), (\ref{Riccidecomp2}) and (\ref{Riccidecomp3}) into the decomposition~(\ref{Riccidecomp}) of the off-block diagonal Ricci curvature components $R_{\jalpha \kbeta}$ to obtain, for $\alpha \neq \beta$,
\begin{gather}
R_{\jalpha \kbeta}=-\frac{1}{2}\sum_{\gamma =1}^{r}\lgamma \partial_{\jalpha}\partial_{\kbeta}\left(\log\left(\frac{\det S}{s^{\gamma 1}}\right)\right)+\partial_{\jalpha}\partial_{\kbeta}\big(\log(\det S)\big)\nonumber\\
\hphantom{R_{\jalpha \kbeta}=}{} + \frac{1}{4}(\lalpha+\lbeta-2) \partial_{\jalpha}\left(\log\left(\frac{\det S}{s^{\beta 1}}\right)\right)\partial_{\kbeta}\left(\log\left(\frac{\det S}{s^{\alpha 1}}\right)\right)\nonumber\\
\hphantom{R_{\jalpha \kbeta}=}{}+\frac{1}{4}\sum_{\gamma \neq \alpha, \,\beta=1}^{r}\lgamma \left[\partial_{\kbeta}\left(\log\left(\frac{\det S}{s^{\alpha 1}}\right)\right) \partial_{\jalpha}\left(\log\left(\frac{\det S}{s^{\gamma 1}}\right)\right)\right.\nonumber\\
\hphantom{R_{\jalpha \kbeta}=}{}
+\partial_{\jalpha}\left(\log\left(\frac{\det S}{s^{\beta 1}}\right)\right)\partial_{\kbeta}\left(\log\left(\frac{\det S}{s^{\gamma 1}}\right)\right)\nonumber\\
\left.\hphantom{R_{\jalpha \kbeta}=}{} -\partial_{\kbeta}\left(\log\left(\frac{\det S}{s^{\gamma 1}}\right)\right)\partial_{\jalpha}\left(\log\left(\frac{\det S}{s^{\gamma 1}}\right)\right)\right].\label{Riccialmostfinal}
\end{gather}
We observe that since the cofactors $s^{\alpha 1}$ and $s^{\beta 1}$ are independent of the groups of variables ${\bf x}^{\alpha}$ and ${\bf x}^{\beta}$ respectively and since $\sum\limits_{\alpha=1}^{r}\lalpha = n$, we have
\begin{gather*}
-\frac{1}{2}\sum_{\gamma =1}^{r}\lgamma \partial_{\jalpha}\partial_{\kbeta}\left(\log\left(\frac{\det S}{s^{\gamma 1}}\right)\right)+\partial_{\jalpha}\partial_{\kbeta}\big(\log(\det S)\big)\\
\qquad{} =-\frac{1}{2}\partial_{j_{\alpha}}\partial_{k_{\beta}}\log\left[\frac{(\det S)^{n-2}}{\big(s^{11}\big)^{l_{1}}\cdots \big(s^{r1}\big)^{l_{r}}}\right].
\end{gather*}
We write
\begin{gather}
-\frac{1}{2}\partial_{j_{\alpha}}\partial_{k_{\beta}}\log\left[\frac{(\det S)^{n-2}}{\big(s^{11}\big)^{l_{1}}\cdots \big(s^{r1}\big)^{l_{r}}}\right]\nonumber\\
\qquad{} =-\frac{3}{4}\partial_{j_{\alpha}}\partial_{k_{\beta}}\log\left[\frac{(\det S)^{n-2}}{\big(s^{11}\big)^{l_{1}}\cdots \big(s^{r1}\big)^{l_{r}}}\right] +\frac{1}{4}\partial_{j_{\alpha}}\partial_{k_{\beta}}\log\left[\frac{(\det S)^{n-2}}{\big(s^{11}\big)^{l_{1}}\cdots \big(s^{r1}\big)^{l_{r}}}\right],\label{simplif}
\end{gather}
and compute the second term in (\ref{simplif}) as follows
\begin{gather*}
\frac{1}{4}\partial_{j_{\alpha}}\partial_{k_{\beta}}\log\left[\frac{(\det S)^{n-2}}{\big(s^{11}\big)^{l_{1}}\cdots \big(s^{r1}\big)^{l_{r}}}\right]\\
\qquad{} =\frac{1}{4}\partial_{j_{\alpha}}\partial_{k_{\beta}}\left[ \sum_{1\leq \gamma \neq \alpha, \beta\leq r} \log\left( \frac{(\det S)^{\lgamma}}{\big(s^{\gamma 1}\big)^{\lgamma}}\right)+\log\left(\frac{(\det S)^{\lalpha+\lbeta-2}}{s^{\alpha 1}s^{\beta 1}}\right)\right].
\end{gather*}
Evaluating the derivatives an using the fact that the cofactors $s^{\alpha 1}$ and $s^{\beta 1}$ are independent of the groups of variables ${\bf x}^{\alpha}$ and ${\bf x}^{\alpha}$, we conclude that the expression (\ref{Riccialmostfinal}) of the off-block diagonal Ricci curvature can be written as
\begin{gather*}
R_{\jalpha \kbeta} =-\frac{3}{4}\partial_{j_{\alpha}}\partial_{k_{\beta}}\log\left[\frac{(\det S)^{n-2}}{\big(s^{11}\big)^{l_{1}}\cdots \big(s^{r1}\big)^{l_{r}}}\right] + \frac{1}{4} (l_{\alpha}+l_{\beta}-2)\frac{s^{\alpha 1}s^{\beta 1}}{\det S}\partial_{j_{\alpha}}\partial_{k_{\beta}}\left(\frac{\det S}{s^{\alpha 1}s^{\beta 1}}\right)\\
\hphantom{R_{\jalpha \kbeta} =}{}+
\sum_{\gamma\neq \alpha,\beta=1}^{r}l_{\gamma}\left(\partial_{j_{\alpha}}\log\frac{s^{\gamma 1}}{\det S}\right)\left(\partial_{k_{\beta}}\log\frac{s^{\alpha 1}}{\det S}\right)
+\left(\partial_{j_{\alpha}}\log\frac{s^{\beta 1}}{\det S}\right)\left(\partial_{k_{\beta}}\log\frac{s^{\gamma 1}}{\det S}\right)\\
\hphantom{R_{\jalpha \kbeta} =}{}
-\frac{\det S}{s^{\gamma 1}}\partial_{\jalpha}\partial_{\kbeta}\left(\frac{s^{\gamma 1}}{\det S}\right).
\end{gather*}
This completes the proof of Theorem \ref{maintheorem1}.
\begin{thebibliography}{99}
\bibitem[1]{Ben2016}
Benenti S., Separability in {R}iemannian manifolds, \href{https://doi.org/10.3842/SIGMA.2016.013}{\textit{SIGMA}} \textbf{12}
 (2016), 013, 21~pages, \href{https://arxiv.org/abs/1512.07833}{arXiv:1512.07833}.

\bibitem[2]{BCR1-2002}
Benenti S., Chanu C., Rastelli G., Remarks on the connection between the
 additive separation of the {H}amilton--{J}acobi equation and the
 multiplicative separation of the {S}chr\"{o}dinger equation. {I}.~{T}he
 completeness and {R}obertson conditions, \href{https://doi.org/10.1063/1.1506180}{\textit{J.~Math. Phys.}} \textbf{43}
 (2002), 5183--5222.

\bibitem[3]{BCR2-2002}
Benenti S., Chanu C., Rastelli G., Remarks on the connection between the
 additive separation of the {H}amilton--{J}acobi equation and the
 multiplicative separation of the {S}chr\"{o}dinger equation. {II}.~{F}irst
 integrals and symmetry operators, \href{https://doi.org/10.1063/1.1506181}{\textit{J.~Math. Phys.}} \textbf{43} (2002),
 5223--5253.

\bibitem[4]{BCR2005}
Benenti S., Chanu C., Rastelli G., Variable-separation theory for the null
 {H}amilton--{J}acobi equation, \href{https://doi.org/10.1063/1.1862325}{\textit{J.~Math. Phys.}} \textbf{46} (2005),
 042901, 29~pages.

\bibitem[7]{Car1977}
Carter B., Killing tensor quantum numbers and conserved currents in curved
 space, \href{https://doi.org/10.1103/PhysRevD.16.3395}{\textit{Phys. Rev.~D}} \textbf{16} (1977), 3395--3414.

\bibitem[8]{CR2006}
Chanu C., Rastelli G., Fixed energy {$R$}-separation for {S}chr\"{o}dinger
 equation, \href{https://doi.org/10.1142/S021988780600120X}{\textit{Int.~J. Geom. Methods Mod. Phys.}} \textbf{3} (2006),
 489--508, \href{https://arxiv.org/abs/nlin.SI/0512033}{arXiv:nlin.SI/0512033}.

\bibitem[9]{CR2007}
Chanu C., Rastelli G., Eigenvalues of {K}illing tensors and separable webs on
 {R}iemannian and pseudo-{R}iemannian manifolds, \href{https://doi.org/10.3842/SIGMA.2007.021}{\textit{SIGMA}} \textbf{3}
 (2007), 021, 21~pages, \href{https://arxiv.org/abs/nlin.SI/0612042}{arXiv:nlin.SI/0612042}.

\bibitem[10]{CR2019}
Chanu C.M., Rastelli G., Block-separation of variables: a form of partial
 separation for natural {H}amiltonians, \href{https://doi.org/10.3842/SIGMA.2019.013}{\textit{SIGMA}} \textbf{15} (2019),
 013, 22~pages, \href{https://arxiv.org/abs/1808.01889}{arXiv:1808.01889}.

\bibitem[11]{DKN2015}
Daud\'e T., Kamran N., Nicoleau F., Inverse scattering at fixed energy on
 asymptotically hyperbolic {L}iouville surfaces, \href{https://doi.org/10.1088/0266-5611/31/12/125009}{\textit{Inverse Problems}}
 \textbf{31} (2015), 125009, 37~pages, \href{https://arxiv.org/abs/1409.6229}{arXiv:1409.6229}.

\bibitem[12]{DKN2016}
Daud\'e T., Kamran N., Nicoleau F., Non-uniqueness results for the
 anisotropic {C}alder\'{o}n problem with data measured on disjoint sets,
 \href{https://doi.org/10.5802/aif.3240}{\textit{Ann. Inst. Fourier (Grenoble)}} \textbf{69} (2019), 119--170,
 \href{https://arxiv.org/abs/1510.06559}{arXiv:1510.06559}.

\bibitem[13]{DKN2017}
Daud\'e T., Kamran N., Nicoleau F., On the hidden mechanism behind
 non-uniqueness for the anisotropic {C}alder\'{o}n problem with data on
 disjoint sets, \href{https://doi.org/10.1007/s00023-018-00755-2}{\textit{Ann. Henri Poincar\'e}} \textbf{20} (2019), 859--887,
 \href{https://arxiv.org/abs/1701.09056}{arXiv:1701.09056}.

\bibitem[14]{DKN2018a}
Daud\'e T., Kamran N., Nicoleau F., The anisotropic {C}alder\'on problem for
 singular metrics of warped product type: the borderline between uniqueness
 and invisibility, \textit{J.~Spectr. Theory}, {t}o appear,
 \href{https://arxiv.org/abs/1805.05627}{arXiv:1805.05627}.

\bibitem[15]{DKN2019a}
Daud\'e T., Kamran N., Nicoleau F., On non-uniqueness for the anisotropic
 {C}alder\'on problem with partial data, \href{https://arxiv.org/abs/1901.10467}{arXiv:1901.10467}.

\bibitem[16]{DKN2019b}
Daud\'e T., Kamran N., Nicoleau F., The anisotropic Calder\'on problem on
 3-dimensional conformally {S}t\"ackel manifolds, \href{https://arxiv.org/abs/1909.01669}{arXiv:1909.01669}.

\bibitem[17]{diPirro1896}
Di~Pirro G., Sugli integrali primi quadratici delle equazioni della mecanica,
 \href{https://doi.org/10.1007/BF02419532}{\textit{Annali di Mat.}} \textbf{24} (1896), 315--334.

\bibitem[18]{Eis1934}
Eisenhart L.P., Separable systems of {S}tackel, \href{https://doi.org/10.2307/1968433}{\textit{Ann. of Math.}}
 \textbf{35} (1934), 284--305.

\bibitem[21]{Gob2016}
Gobin D., Inverse scattering at fixed energy on three-dimensional
 asymptotically hyperbolic {S}t\"{a}ckel manifolds, \href{https://doi.org/10.4171/PRIMS/54-2-2}{\textit{Publ. Res. Inst.
 Math. Sci.}} \textbf{54} (2018), 245--316, \href{https://arxiv.org/abs/1605.05115}{arXiv:1605.05115}.

\bibitem[22]{GT2013}
Guillarmou C., Tzou L., The {C}alder\'{o}n inverse problem in two dimensions,
 in Inverse Problems and Applications: Inside Out.~{II}, \textit{Math. Sci.
 Res. Inst. Publ.}, Vol.~60, Cambridge University Press, Cambridge, 2013,
 119--166.

\bibitem[23]{HaMa1978}
Hauser I., Malhiot R.J., Forms of all spacetime metrics which admit {$[(11)
 (11)]$} {K}illing tensors with nonconstant eigenvalues, \href{https://doi.org/10.1063/1.523523}{\textit{J.~Math.
 Phys.}} \textbf{19} (1978), 187--194.

\bibitem[25]{KaMi1980}
Kalnins E.G., Miller Jr. W., Killing tensors and variable separation for
 {H}amilton--{J}acobi and {H}elmholtz equations, \href{https://doi.org/10.1137/0511089}{\textit{SIAM~J. Math. Anal.}}
 \textbf{11} (1980), 1011--1026.

\bibitem[26]{KaMi1981}
Kalnins E.G., Miller Jr. W., Killing tensors and nonorthogonal variable
 separation for {H}amilton--{J}acobi equations, \href{https://doi.org/10.1137/0512054}{\textit{SIAM~J. Math. Anal.}}
 \textbf{12} (1981), 617--629.

\bibitem[27]{KaMi1984}
Kalnins E.G., Miller Jr. W., The theory of orthogonal {$R$}-separation for
 {H}elmholtz equations, \href{https://doi.org/10.1016/0001-8708(84)90004-5}{\textit{Adv. Math.}} \textbf{51} (1984), 91--106.

\bibitem[28]{KKL2001}
Katchalov A., Kurylev Y., Lassas M., Inverse boundary spectral problems,
\textit{Chapman \& Hall/CRC Monographs and Surveys in Pure and Applied
 Mathematics}, Vol.~123, \href{https://doi.org/10.1201/9781420036220}{Chapman \& Hall/CRC}, Boca Raton, FL, 2001.

\bibitem[29]{KS2014}
Kenig C., Salo M., Recent progress in the {C}alder\'{o}n problem with partial
 data, in Inverse Problems and Applications, \href{https://doi.org/10.1090/conm/615/12245}{\textit{Contemp. Math.}}, Vol.~615, Amer. Math. Soc., Providence, RI, 2014, 193--222, \href{https://arxiv.org/abs/1302.4218}{arXiv:1302.4218}.

\bibitem[32]{LU2001}
Lassas M., Uhlmann G., On determining a {R}iemannian manifold from the
 {D}irichlet-to-{N}eumann map, \href{https://doi.org/10.1016/S0012-9593(01)01076-X}{\textit{Ann. Sci. \'Ecole Norm. Sup.~(4)}}
 \textbf{34} (2001), 771--787.

\bibitem[33]{LeU1989}
Lee J.M., Uhlmann G., Determining anisotropic real-analytic conductivities by
 boundary measurements, \href{https://doi.org/10.1002/cpa.3160420804}{\textit{Comm. Pure Appl. Math.}} \textbf{42} (1989),
 1097--1112.

\bibitem[36]{Mi1988}
Miller Jr. W., Mechanisms for variable separation in partial differential
 equations and their relationship to group theory, in Symmetries and Nonlinear
 Phenomena ({P}aipa, 1988), \textit{CIF Ser.}, Vol.~9, World Sci. Publ.,
 Teaneck, NJ, 1988, 188--221.

\bibitem[37]{Pain1897}
Painlev\'e P., Sur les int{\'e}grales quadratiques des \'equations de la
 dynamique, \textit{C.~R.~Acad. Sci. Paris} \textbf{124} (1897), 221--224.

\bibitem[38]{Per1990}
Perelomov A.M., Integrable systems of classical mechanics and {L}ie algebras~
 {V}ol.~{I}, \href{https://doi.org/10.1007/978-3-0348-9257-5}{Birkh\"{a}user Verlag}, Basel, 1990.

\bibitem[39]{Rob1928}
Robertson H.P., Bemerkung \"{u}ber separierbare {S}ysteme in der
 {W}ellenmechanik, \href{https://doi.org/10.1007/BF01451624}{\textit{Math. Ann.}} \textbf{98} (1928), 749--752.

\bibitem[40]{Sa2013}
Salo M., The {C}alder\'{o}n problem on {R}iemannian manifolds, in Inverse
 problems and applications: inside out.~{II}, \textit{Math. Sci. Res. Inst.
 Publ.}, Vol.~60, Cambridge University Press, Cambridge, 2013, 167--247.

\bibitem[41]{Sta1897}
St\"ackel P., Sur les probl\`emes de la dynamique, qui se r\'eduisent a des
 quadratures, \textit{C.~R.~Acad. Sci. Paris} \textbf{124} (1897), 221--224.

\bibitem[42]{Tay2011a}
Taylor M.E., Partial differential equations. {I}.~{B}asic theory,
 \textit{Applied Mathematical Sciences}, Vol.~115, \href{https://doi.org/10.1007/978-1-4684-9320-7}{Springer-Verlag}, New York,
 1996.

\bibitem[43]{Top1999}
Topalov P., Integrability criterion of geodesical equivalence. {H}ierarchies,
 \href{https://doi.org/10.1023/A:1006369525091}{\textit{Acta Appl. Math.}} \textbf{59} (1999), 271--298.

\bibitem[44]{Top2002}
Topalov P., Commutative conservation laws for geodesic flows of metrics
 admitting projective symmetry, \href{https://doi.org/10.4310/MRL.2002.v9.n1.a5}{\textit{Math. Res. Lett.}} \textbf{9} (2002),
 65--72.

\bibitem[46]{Uhl2009}
Uhlmann G., Electrical impedance tomography and {C}alder\'{o}n's problem,
 \href{https://doi.org/10.1088/0266-5611/25/12/123011}{\textit{Inverse Problems}} \textbf{25} (2009), 123011, 39~pages.

\bibitem[47]{Uhl2014}
Uhlmann G., Inverse problems: seeing the unseen, \href{https://doi.org/10.1007/s13373-014-0051-9}{\textit{Bull. Math. Sci.}}
 \textbf{4} (2014), 209--279.

\end{thebibliography}
\end{document}
